\documentclass[fontset=fandol, a4paper, 12pt, twoside]{ctexart}

% ==========================================
% 宏包加载
% ==========================================
\usepackage{geometry}       % 页面布局
\usepackage{hyperref}       % 超链接
\usepackage{fancyhdr}       % 页眉页脚
\usepackage{titlesec}       % 标题格式定制
\usepackage{enumitem}       % 列表定制
\usepackage{setspace}       % 行距
\usepackage{xcolor}         % 颜色
\usepackage{footmisc}       % 脚注定制

% ==========================================
% 页面布局设置 (双面批注模式)
% ==========================================
% twoside 模式下：
% inner = 装订侧 (奇数页左, 偶数页右)
% outer = 翻阅侧 (奇数页右, 偶数页左) -> 这里留出 6.5cm 给您做笔记
\geometry{
    a4paper,
    inner=2.5cm,     
    outer=6.5cm,     
    top=3cm, 
    bottom=3cm,
    marginparwidth=5cm,
    marginparsep=0.5cm
}

% 行距设置：1.5倍行距，方便阅读和行间标注
\linespread{1.5}

% ==========================================
% 样式定制
% ==========================================

% 1. 链接颜色设置
\hypersetup{
    colorlinks=true,
    linkcolor=black,    
    urlcolor=blue,      
    citecolor=blue,
    pdftitle={矛盾理论},
    % pdfsubject={ISBN: },
    pdfauthor={阿兰·巴迪欧}
}

% 2. 页眉页脚设置 (双面模式)
\pagestyle{fancy}
\fancyhf{} % 清空默认设置

% 奇数页 (Odd): 页码在右(外), 标题在左(内)
\fancyhead[RO]{\small \thepage} 
\fancyhead[LO]{\small \textit{\leftmark}} 

% 偶数页 (Even): 页码在左(外), 标题在右(内)
\fancyhead[LE]{\small \thepage} 
\fancyhead[RE]{\small \textit{矛盾理论 / 阿兰·巴迪欧}} 

\renewcommand{\headrulewidth}{0.5pt}

% 3. 标题格式定制
\ctexset{
    section = {
        format = \Large\bfseries,
        number = \arabic{section},
        name = {, .},
        aftername = \hspace{0.5em},
        beforeskip = 20pt,
        afterskip = 15pt
    },
    subsection = {
        format = \large\bfseries,
        number = \Alph{subsection},
        name = {, .},
        aftername = \hspace{0.5em},
        beforeskip = 15pt,
        afterskip = 10pt
    },
    subsubsection = {
        format = \normalsize\bfseries,
        number = {},
        beforeskip = 10pt,
        afterskip = 5pt
    }
}

% 4. 引用环境定制
\renewenvironment{quote}
  {\list{}{\rightmargin=0cm \leftmargin=2em}%
   \item\relax\small\kaishu\color{black}}
  {\endlist}

% ==========================================
% 文档正文
% ==========================================
\begin{document}

% % --- 封面 ---
% \begin{titlepage}
%     \centering
%     \vspace*{3cm}
%     {\Huge \bfseries 矛盾理论}\\[1.5cm]
%     {\LARGE \textbf{Théorie de la contradiction}}\\[2cm]
%     {\Large \textbf{【法】阿兰·巴迪欧}}\\[0.5cm]
%     {\large 1975年5月}\\[3cm]
    
%     \begin{quote}
%     \centering
%     \textit{“马克思主义的道理千条万绪，归根结底，就是一句话：造反有理。”}
%     \end{quote}
% \end{titlepage}

% \newpage
% % 目录 (可选，如果不需要可注释掉)
% \tableofcontents
% \newpage

% --- 正文开始 ---
\setcounter{page}{1}

\section{一个重要的哲学论点：“造反有理”}

毛泽东有一句著名的论断：“马克思主义的道理千条万绪，归根结底，就是一句话：造反有理。”

这句话如此简单，同时又相当神秘：马克思那庞大的理论事业，他那经过千锤百炼、毫不疏忽的分析，怎么可能全部浓缩在一句格言中：“造反有理”？这句格言又是什么？

这是一种陈述，总结了马克思主义对客观矛盾的分析，以及革命与反革命之间不可避免的对抗吗？这是一种旨在动员革命力量的主观指令吗？马克思主义的真理是：人们反抗了，所以是正确的？还是更确切地说：必须反抗？也许两者兼而有之，更重要的是从一个到另一个的螺旋式运动，即现实的反抗（客观力量）在对其正确性（主观力量）的认识中得到丰富和回归。

\subsubsection*{实践、理论、认识}

在这里，我们已经获得了某种本质性的东西：任何马克思主义的论断，都是一种分裂为二的统一运动，既是陈述又是指令。它浓缩了现实的实践，与实践的运动相一致，并回归于实践。既然存在的事物只在它的发展中才有其存在，那么作为理论的存在——对存在事物的认识——也只在向其所论述的事物运动中才有其存在。任何认识都是一种导向，任何描述都是一种规定。

“造反有理”这句话比其他任何话都更能证明这一点。它表明，马克思主义在成为一门发达的社会形态科学之前，是反抗所要求的缩影：即赋予其正确性。马克思主义是一种党派立场，也是一种党派经验的系统化。社会形态科学的存在之所以对群众有意义，是因为它反映并浓缩了他们的真实革命运动。

马克思主义必须被看作是人民革命累积的智慧，是革命所产生的正确性，是对革命目标的固定和精确化。毛泽东的这句话清楚地将反抗定位为正确思想的起源地，将反动派定位为理论上应当被摧毁的对象。毛泽东的这句话将马克思主义的真理定位在理论与实践的统一之内。

马克思主义的真理是反抗为打倒敌人而形成的正确性。它摒弃了真理面前人人平等的观念。在描述和指令这两种特定的分裂运动中，它对敌人进行评判，宣判其罪行，并沉浸在执行判决的行动中。

反抗者们根据他们已有的本质运动，认识到他们的权力和义务：消灭反动派。马克思的《资本论》说的正是这一点：无产阶级有理由暴力推翻资本家。马克思主义的真理不是一种调和的真理。它本身就是专政，必要时甚至是恐怖。

毛泽东的这句话提醒我们，对于一个马克思主义者来说，理论与实践（正确性与反抗）的结合是理论本身的内在条件，因为真理是现实的过程，真理就是反抗反动派。黑格尔的这段论述再深刻、再真实不过了：

\begin{quote}
“绝对理念，如其所示，是实践理念和理论理念的同一性，两者各自都还是片面的。”\footnote{黑格尔，《逻辑学》，奥比尔出版社，第二卷，第549页。}
\end{quote}

对于黑格尔来说，绝对真理是理论与实践的矛盾统一。它是存在与行动之间不间断且分裂的过程。列宁对此给予了热情的赞扬：

\begin{quote}
“理论理念（认识）和实践的统一——这很重要——而这种统一正是在认识论中实现的，因为其结果就是绝对理念。”\footnote{列宁，《哲学笔记》，社会出版社，第208页。}
\end{quote}

让我们非常仔细地阅读这句话，因为引人注目的是，它把“认识”这个词分成了两部分。这是一个关键点，我们稍后会经常回到这一点：认识，作为理论，与实践（辩证地）对立。理论与实践构成一个统一体，也就是说，对于辩证法来说，是两个对立面的统一。

但是，认识（理论）/实践这个矛盾，反过来又是认识论本身的客体。换句话说：认识过程的内在性质是理论/实践的矛盾。或者说：实践，就其本身而言，辩证地与认识（理论）对立，但仍然是认识作为过程的一个组成部分。

在所有的马克思主义文献中，我们都能发现这种分裂，这种“认识”一词的双重出现，这取决于它指的是与实践辩证相关的理论，还是指这种辩证的整体过程，即理论与实践这两个术语的矛盾运动。

看看毛泽东的《人的正确思想是从哪里来的？》：

\begin{quote}
“要完成从实践到认识、又从认识到实践的多次反复，才能完成正确的认识。这就是马克思主义的认识论，辩证唯物主义的认识\mbox{论。”\footnote{毛泽东，《人的正确思想是从哪里来的？》，《毛泽东著作选集》第五卷，人民出版社，1977年，第309页。}}
\end{quote}

认识的运动，就是实践-认识-实践的交替。这里的“认识”既指过程中的一个环节，也指整个过程，而这个过程反过来又包括了最初和最终的两个实践环节。

为了固定语言\footnote{马克思列宁主义毛泽东思想不是形式主义。它所用的词语都融入了现实认识的破坏/建设运动中。如果目标达成，符号并不重要。因此，词语可以移动：重要的是它们的力量。在这里，力量也胜过了对位置的尊重。}（以及传统），我们将把理论称为理论/实践矛盾的一个环节，而其整体运动，则将是认识的过程。我们将说：认识是实践/理论的辩证过程。

由此可见，那些幻想能够绕过实践第一性这一战略论点的人，他们的反动幻觉是多么明显。很清楚，任何不处于现实革命运动中，任何不在实践上内在于反抗反动派的人，即使他进行理论化，也一无所知。

诚然，毛泽东曾指出，在理论/实践的矛盾中，即在现实过程的一个阶段中，理论可以暂时起主要作用：

\begin{quote}
“当列宁所说的‘没有革命的理论，就不会有革命的运动’的那种情形出现的时候，革命理论的创造和提倡，就起了主要地、决定地作用。”\footnote{毛泽东，《矛盾论》，《毛泽东著作选集》第一卷，人民出版社，1991年，第323-324页。}
\end{quote}

这是否意味着理论在此时是一种内在的革命可能性，以至于“纯粹的马克思主义理论家”可以而且应该涌现出来？绝非如此。

它意味着，在作为认识过程的理论/实践矛盾中，理论是矛盾的主要方面；是革命实践经验的系统化使得前进成为可能；继续定量地累积和重复这些经验是毫无意义的，因为此时的议程是质的飞跃，是紧随其应用，即紧随其验证而来的理性综合。

但是，没有这些经验，没有有组织的实践（因为只有组织才能集中经验），就不会有系统化，根本就不会有认识。没有有组织的应用，就不会有验证计划，就不会有验证，就不会有真理。此时的“理论”只能产生唯心主义的荒谬。

因此，我们回到起点：实践是真理理性运动的内在部分。在与理论的对立中，它是认识的一部分。正是这种直觉让列宁对黑格尔关于绝对理念的构想充满热情，以至于他将马克思视为黑格尔的简单继承者。（“因此，马克思直接地继承了黑格尔，他把实践标准引入了认识论。”\footnote{列宁，《哲学笔记》，社会出版社，第201页。}）

毛泽东的这句话精确了列宁的热情。它是黑格尔辩证论断的一般历史内容。并非任何实践都是理论的内在锚点，而是反抗反动派的实践。而理论反过来，并非外在于实践、外在于反抗，对其实施管辖：它通过对其正确性的中介性阐明而融入其中。从这个意义上说，这句话确实说尽了一切，它概括了马克思主义的阶级立场，其具体的革命意义。

任何试图从决裂而非反抗的角度；从体系而非历史的角度；从理论第一性而非实践第一性的角度；从无产阶级作为先验条件而非工人阶级累积的智慧的角度来审视马克思主义的人，都置身于这个整体之外。

\subsubsection*{“正确性”（raison）一词的三种含义}

然而，如果这句话说尽了一切，那它也是按照辩证法的方式，即以一种分裂的简单性来表达的。集中并表面上掩盖这种分裂的，是“正确性”（raison）这个词：人们有正确性，反抗有正确性，一种新的正确性正在对抗反动派。事实上，通过“正确性”这个词，这句话说了三件事，而正是这三者的结合构成了整体。

A. \textbf{造反有理，这首先意味着，并非“必须反抗反动派”，而是：人们反抗反动派，这是一个事实，这个事实就是正确性。}这句话表明了：实践第一性。

反抗不需要等待其正确性，反抗是无论出于任何可能的原因而总是早已存在的。马克思主义只是说：反抗就是正确性，反抗就是主体。马克思主义是反抗智慧的总结。为什么要去写《资本论》，写成百上千页的细致入微、刻苦钻研的智慧，写那些有时难以理解的辩证法著作？因为只有这样才能与反抗的深刻智慧相称。

反抗的历史深度和顽强性先于马克思主义而存在，它累积了其出现的条件和必然性，因为它根植了一种信念：超越引发无产阶级起义的特定原因，存在一种深刻且无法根除的正确性。马克思的《资本论》正是将历史上所有原因的总和中所呈现出的东西，以一般正确性的方式进行系统化。

资产阶级虽然了解并承认阶级斗争，也愿意承认和寻找反抗的特定原因，哪怕只是为了防止其再次发生。但它却忽视了无产阶级最终所秉持的正确性，这种正确性是任何消除原因和环境的措施都无法满足的。马克思的事业是反映出，不仅仅是在具体的斗争中，更是在其中所投入的阶级能量的持久性和发展中，所呈现出的东西。仅仅对原因的思考是不够的。

（列宁在为黑格尔辩护、反对康德时，很好地强调了因果性范畴的不足之处，因为黑格尔没有给它一个特殊的地位：“当人们读黑格尔关于因果性的论述时，乍一看似乎很奇怪，他为何如此轻描淡写这个康德派如此珍爱的主题。那是因为对他来说，因果性是普遍联系的众多规定性之一……”\footnote{列宁，《哲学笔记》，社会出版社，第152页。}）

我们必须从深层次上解释这种持久性。无产阶级立场的本质不在于阶级斗争的插曲，而在于其双后所隐藏的历史计划，这个计划的不知疲倦的持久性和无产阶级顽强性的连续阶段，是其实践存在的形式。正确性就存在于那里。

只有对它的澄清和阐述，同时作为反映和指令，才能真正公正地对待反抗所揭示的现象的阶级存在运动。

今天，只有毛泽东的事业才充分发展了无产阶级在其无条件和持久的反抗中所做和所揭示的一切。只有它才说了：是的，矛盾是对抗性的，是的，作为这种矛盾之火的工人反抗，正是历史的正确性。造反有理，首先意味着：顽强的无产阶级是正确的，所有的正确性都在他们一边，甚至更多。

B. \textbf{“造反有理”，也意味着：反抗将会胜利。} 反动派将在历史的法庭上为所有的剥削和压迫罪行付出代价。

无产阶级反抗的顽强性，当然是“正确性”一词的第一层含义，即工人与资产阶级对立的矛盾的客观的、不可化解的性质，但它也是最终胜利的实践确定性，是对工人失败主义的自发的、不断重生的批判。事物的现状是不可接受的，是分裂的，这是反抗反动派的第一个正确性。

现状是暂时的，是注定要失败的，这是第二个正确性。这是从未来的角度，而不是从动机或当下的角度来看的正确性。这是胜利意义上的正确性，超越了合法性意义上的正确性。反抗是智慧，因为它公正，因为它基于正确性，也因为它决定着未来。马克思主义摒弃任何仅仅是辩护性的正确性观念。无产阶级不仅有反抗的真实理由，更有胜利的理由。

“正确性”在这里是革命合法性与革命乐观主义的交汇点。
反抗对康德的道德格言“你必须，所以你能够”是过敏的。康德得出结论说，这样纯粹由责任决定的行为很可能从未发生过。道德是一种失败的规定。但工人反抗确实发生了，并在马克思主义中找到了其胜利的规定之地。马克思主义的正确性不是应然，而是对存在本身的肯定，是对那种起义、反对、矛盾的事物的无限力量的肯定。它是人民拒绝的客观胜利。作为唯物主义的，工人的正确性是说：“你能够，所以你必须。”

C. \textbf{“造反有理”这次意味着：反抗可以通过对其自身正确性的认识而得到加强。} 这句话本身“造反有理”既是反抗本身内在的认识内核的发展，也是这种发展回归到反抗中去。反抗，本身具有正确性，在马克思主义中找到了发展这种正确性、确保其胜利正确性的方法。

使得反抗的合法性（“正确性”一词的第一种含义）与其胜利（“正确性”一词的第二种含义）相互关联的，是一种反抗作为始终存在的实践与其正确性的发达形式之间的新型融合。马克思主义与现实工人运动的融合，这就是“正确性”一词的第三种含义，即其前两种含义的辩证的、客观的和主观的联系。

我们在这里再次找到了马克思主义论断的辩证地位，它们总是根据反映和指令而分裂：通过超越原因抓住阶级能量的正确性，理论同时阐明了正确性可以胜过原因、整体可以胜过局部、战略可以胜过战术的规则。反抗在其实践的持久性中阐述其正确性；但对这种正确性的澄清表述打破了这种持久性的重复规则。

反抗以其自身的正确性武装自己，而不仅仅是展现它。它集中了其理性的特质：它组织其正确性，并部署其胜利的工具。

认识到造反有理，通过揭示这种实践正确性背后的理论正确性，使得主观（组织、计划）能够与客观（阶级斗争、反抗）相匹敌。

“正确性”曾代表革命的合法性和乐观主义，现在则代表了对历史的自觉和掌控。

\subsubsection*{作为矛盾的正确性}

“造反有理”这句话确实说尽了历史运动的一切，因为它说出了运动的能量、方向和工具。能量是阶级斗争，是反抗内在的客观理性。方向是剥削和压迫世界的必然崩溃，是共产主义的正确性。

工具是历史中能量与方向之间可能的关系指导，即阶级斗争（始终且处处是历史的动力）与共产主义计划（始终且处处是受压迫者反抗所高举的价值）之间的关系指导。

这个工具就是成为主体的正确性，就是党。

“造反有理”说尽了整体，因为它说出了阶级斗争和实践第一性，共产主义和国家的消亡，党和无产阶级专政。

这句话说的是完整的正确性，即按照主观和客观、现实和计划、目标和阶段而分裂的正确性。

而这种完整的正确性就是矛盾：无法独自拥有正确性，无法为自己而拥有正确性。造反有理，是反抗反动派。我们总是反抗反动派，这句“反抗反动派”是真理的内在条件。

这就是为什么毛泽东的这句话总结了马克思主义；它表明：任何正确性都包含矛盾。“正确思想是在同错误思想的斗争中形成的”，正确性是在反抗不正确性、反抗中国人一贯称之为“反动谬论”的斗争中锻造的。

一切真理都在对无意义的摧毁中得到肯定。因此，一切真理本质上都是摧毁。一切只求保持不变的事物都只是一种虚假。马克思主义的认识领域永远是一个废墟。

毛泽东的这句话向我们揭示了辩证法：作为反抗的正确性的阶级本质存在于对立面的你死我活的斗争中。真理只存在于一个分裂的过程中。

矛盾论完全包含在反抗者的历史智慧中。因此，辩证法和反抗一样，自古以来就存在。辩证法从哲学上概括了反抗现有世界并渴望其彻底改变的被剥削者的世界观。这就是为什么它是一种永恒的哲学倾向，它不懈地反对保守的形而上学压迫：

\begin{quote}
“在人类认识史上，关于世界发展的观点，从来就有两种：一种是形而上学的观点，一种是辩证法的观点，形成了互相对立的两种宇宙观。”\footnote{毛泽东，《矛盾论》，《毛泽东著作选集》第一卷，人民出版社，1991年，第323-324页。}
\end{quote}

我们总是要继续辩证法，使其对抗形而上学，这意味着：赋予反抗者正确性。今天，就是赋予真正的马克思主义以正确性，对抗虚假的马克思主义。

赋予毛主义者正确性，对抗修正主义者。

\section{探寻辩证法的基本原则}

赋予反抗者正确性，反抗反动派，这在哲学上意味着：在历史过程的每一个阶段，重新统一被剥削者的思想：即辩证法思想；使这种统一的原则变得生动和具有战斗性。

但是，辩证法是否有既定的原则？有多少？是哪些？所有无产阶级领导人都讨论过这个问题。我们将看到，他们并没有就此画上句号。正如辩证法思想并不开始于反抗的运动中，它也永远不会结束，其最先进的系统化表述，也永远只是无产阶级领导人的临时哲学指令。

\subsection{恩格斯}

所有革命思想都部分地由其所发展的对象决定。恩格斯和马克思所进行的自觉工人先锋队的第一场重大思想斗争，主要是针对各种小资产阶级唯心主义：空想社会主义；蒲鲁东的联邦主义；巴枯宁的无政府主义。

为了摧毁所有的急躁和幻想，他们要做的就是恢复现实革命运动的客观性，恢复阶级斗争的物质基础。因此，在规定无产阶级哲学的矛盾中——唯物主义/辩证法——恩格斯最常倾向于将唯物主义论点视为主要方面。

然而，在那个时代，在继承革命资产阶级的思想斗争中，任何唯物主义都必须以自然科学为依据，因为自然科学仍然是摧毁宗教和唯心主义构建的有力工具。面对“空想社会主义”，恩格斯自然而然地以“科学社会主义”来对抗。

恩格斯无疑是一位伟大的辩证法家。看看他具体的历史分析吧。

然而，在他对辩证法的总体（哲学）阐述中，他将其严格地从属于唯物主义，并最终从属于自然科学的内容。这并非一种偏差，而是工人运动中两条道路斗争的历史性、本质性要求。

对于恩格斯来说，科学对辩证法的吸收是主要目标，而这一过程最终将哲学消解在科学本身中，只留下一个认识论的残余：

\begin{quote}
“只有当自然科学和历史科学吸收了辩证法时，所有哲学废物——纯粹的思想理论除外——都将变得多余，并消失在实证科学中。”\footnote{恩格斯，《自然辩证法》，社会出版社，1961年，第211页。}
\end{quote}

（在《费尔巴哈》中也有同样的主题，其结果是确定了哲学的现实领域：智力过程：“从此以后，哲学只剩下……关于思想过程本身的规律的学说，即逻辑学和辩证法。”\footnote{恩格斯，《路德维希·费尔巴哈与德国古典哲学的终结》，作品全集，第三卷，第397页。} 奇怪的是：这个唯物主义的结果听起来有点唯心主义。）

这种哲学的论战性规定所带来的主要后果是，在恩格斯那里，矛盾论不是辩证法的重心。他给出的定义在这方面很有特点：为了使其与形而上学对立，他将其定义为“联系的科学”\footnote{恩格斯，《反杜林论》，作品全集，第三卷，第488页。}。

在《路德维希·费尔巴哈》中，他将辩证法称为“一般运动规律的科学”。看来恩格斯将联系和相互依赖的理论置于矛盾理论之上。他试图证明辩证法也属于唯物主义论点，并主张，即使是与黑格尔对立，辩证法规律也并非属于理念的运动，而只是客观规律的简单反映，这些规律在自然科学和历史科学中也同样存在。

这种唯物主义的修正只能导致对辩证法基本规律的一种尚且局限的理解。在哲学与反抗之间，有一个在哲学本身中就存在的联系，那就是唯物主义论点和辩证法论点的不平衡发展。当工人运动中的主要敌人不再是小资产阶级“极左主义”，而是修正主义经济主义时，辩证法论点也将迎来恩格斯只赋予唯物主义论点的充分发展。

因此，对我们来说，在《自然辩证法》中对辩证法规律的表述显得不足。只有我们所处的历史观察点才能进行这种评判。

这些规律是什么？我们知道恩格斯在一篇非常著名的文章中提出了三条：

\begin{quote}
“因此，辩证法的规律是从自然界和人类社会的历史中抽象出来的。它们恰恰是这两个历史发展阶段以及思想本身的最一般规律。它们主要归结为以下三条规律：
\begin{enumerate}
    \item 量变质变的规律；
    \item 对立面相互渗透的规律；
    \item 否定之否定的规律。”\footnote{恩格斯，《自然辩证法》，社会出版社，第69页。}
\end{enumerate}
\end{quote}

乍一看，最引人注目的是作为辩证法原则的纯粹的缺失：即矛盾对同一性的第一性。

由于这一缺失，三条规律都无法得到充分发展：它们仍然是粗略的。认识的渐进性、接近性特征，压倒了反映的特征：这是辩证法论点严格从属于唯物主义论点的痕迹。

将量变质变的规律放在首位，是从一个很小的角度来探讨辩证法问题。

首先，因为正如列宁和毛泽东所指出的，这条“规律”实际上只是唯一真正原则的一个特例：对立面统一，以及每个对立面向另一个对立面的转化。

其次，因为这条规律的辩证本质并非仅仅是肯定转化的可逆性，而是质的飞跃跟随量的积累而来的概念。这一原则的真正含义在于通过飞跃来构想发展，从而在于辩证过程的阶段性理论。在对其规律的简单陈述中，恩格斯忽略了使其成为反对进化论和线性发展观念的武器的东西。他将质的飞跃、阶段和时期的理论，简化为一种简单的转化规律。

毛泽东据说在一篇被斯图亚特·施拉姆（Stuart Schram）\footnote{斯图亚特·施拉姆，《毛泽东未经排练的谈话》，企鹅图书出版社，1974年。}归于他的文章中说过，他认为“否定之否定”根本就不存在。

事实上，我们可以勉强地说无产阶级革命是对资产阶级社会的否定，尽管在无产阶级专政下阶级斗争仍在继续，但我们更应该强调的是资产阶级社会内部的分裂过程、无产阶级与资产阶级矛盾的不间断的阶段性发展。否定这一概念本身，因其带有强烈的唯心主义（黑格尔主义）辩证法色彩，而只是一种方便的缩写。

但是，例如，说无产阶级革命是对封建主义的否定之否定，是完全无法理解的，特别是如果我们只承认“否定之否定”的特定含义，即黑格尔的含义：那将等同于宣称，通过资产阶级否定性的中介，无产阶级产生了封建秩序的完整概念，无产阶级专政是这种秩序的自在自为的本质！

如果我们仔细审视马克思对“否定之否定”范畴的使用，我们会发现这仅仅是指在一个特定阶段，为必然过渡到下一阶段而积累客观条件的一个矛盾过程。

换句话说：一个矛盾运动所摧毁的东西，为摧毁者本身的分裂做好了准备。看看这段有代表性的引文：

\begin{quote}
“资本主义的占有方式，是与资本主义的生产方式相适应的，它对那种以自己劳动为基础的个人私有制来说，是第一个否定。但是，资本主义生产又以自然过程的必然性，造成了对自身的否定。这是否定之否定。它不是重新建立劳动者的私有制，而是建立在资本主义时代的成就之上，即建立在协作和对生产资料包括土地的公共占有之上的个人所有制。”\footnote{马克思，《资本论》，作品全集，第二卷，第152页。}
\end{quote}

“否定之否定”仅仅意味着，通过资本的集中，资产阶级在消灭了独立劳动者的私有财产的同时，使得无产阶级消灭所有私有财产成为必要和客观上可能。

这是因为私有财产不再分散在大量劳动者手中，而是集中在少数资本家手中。换句话说：资本主义积累与小私有财产之间的矛盾为新的矛盾——无产阶级与资本主义财产本身之间的矛盾——扫清了道路。“否定之否定”不是一条规律。它只是一个缩写，用来指代两个矛盾过程的历史衔接。

因此，对立面相互渗透的规律是唯一真正的辩证法规律。然而，恩格斯的表述更强调相互作用，而不是对抗。

人们甚至会怀疑，相互作用规律的巨大扩展是否会使人们忽略了过程的动力、其真正的内在原因，即矛盾。在恩格斯那里，有一种顽强且正当的意图，就是要彻底终结因果原则的一切形而上学表述。事实上，将原因和结果作为封闭实体分开，会妨碍人们抓住互换和反向作用的现象。

第二国际的经济主义，将上层建筑（特别是国家）中的所有革命行动都从属于客观经济条件的“成熟”，这要求进行一场反对反辩证的决定论观念的哲学斗争。原因和结果的割裂及不可改变的顺序，是粗鄙的机会主义的“理论”反映。

然而，恩格斯可能走得有点远了：

\begin{quote}
“所有这些人所缺乏的，就是辩证法。他们只看到这里是原因，那里是结果。他们没有看到，这是一种贫乏的抽象，这种形而上学的、两极对立的东西在现实世界中只在危机期间存在；整个巨大的发展是在相互作用的形式下进行的（尽管力量不平衡，经济运动是最强大、最原始、最决定性的）；这里没有什么绝对的东西，一切都是相对的。对他们来说，黑格尔从未存在过。”\footnote{恩格斯，《致康拉德·施密特信》，1890年10月，作品全集，第三卷，第520页。}
\end{quote}

当一种科学主义的修正主义出现时，援引黑格尔和辩证法总是必不可少的。但话说回来，将“两极对立”限制在危机时刻，只为矛盾思想留下了一个狭窄的领域。

可以说，在这里，辩证法的定义留下了在对立面统一中，在相互作用的掩护下，一种调和性地偏爱统一胜过对立的风险。此外，与恩格斯所说的相反，确实存在绝对的东西，列宁和毛泽东都将有力地指出这一点。恰恰是斗争是绝对的，只有统一是相对的。

最终，恩格斯的三条“规律”削弱了矛盾原则，正是以这种削弱为代价，恩格斯才能主要依赖自然科学来阐明他的论战。因此，在其他历史环境下，一整批教条主义者以这三条“规律”为食。

阶级斗争在继续，辩证法也在继续。今天的恩格斯对我们来说还不够。恩格斯必须得到发展。

\subsection{斯大林}

斯大林是第一个无产阶级专政国家的领导人。他对这个国家现实运动的看法，在很大程度上是基于集中和加速发展生产力，以重工业为核心的要求。为此，必须打破所有障碍，必要时，并越来越多地通过恐怖手段。必须进行一场自觉的、国家性的、暴力的阶级斗争。

斯大林关于辩证法问题的主要论述，见于《辩证唯物主义与历史唯物主义》一文。这篇文章写于1938年：在对党内反对派进行恐怖清洗之后，第二次世界大战前夕。

同样，环境、路线选择和对手都带来了局限性：对工业积累的重视，使得作为质的飞跃承诺的量变得到优先；帝国主义的包围，对外国颠覆的痴迷，反映在所有现象的相互依赖中。斯大林所特有的巨大希望和暴力意志主义，体现在他坚信发展中的事物必然会战胜一切。

所有这一切再一次使得矛盾论带上了一种僵硬的、分明的色彩，以及一种潜在的对进化唯物主义的偏爱，即对我们所称的动态唯物主义的偏爱，而非内部裂变和冲突序列发展的理论。

斯大林自己提出了四个原则，他称之为“马克思主义辩证方法的基本特征”：

\begin{enumerate}
    \item 一切现象普遍相互依赖。这可以称为总体性规律。
    \item 运动原则：“根据辩证法，只有正在诞生和发展的事物才是不可战胜的。”
    \item 量变到质变的转化，这次明确地被构想为量变渐进积累与突然的质变飞跃之间的辩证法，使得事物从一个状态进入另一个状态。
    \item 矛盾规律，被构想为“发展过程的内在内容”。\footnote{斯大林，《辩证唯物主义与历史唯物主义》，《列宁主义问题》，诺尔曼·白求恩出版社，第二卷，第785页。}
\end{enumerate}

斯大林所阐述的四条规律，对我们来说（即从我们的政治需求来看），无疑是对恩格斯三条规律的显著改进。我们满意地看到“否定之否定”消失了。量变质变的转化规律表述得更加严谨；矛盾原则被很好地确立为过程的内在规律。

尽管如此，这四条规律构成了一个异质的集合，同样没有明确地从属于矛盾对同一性的第一性这一论断。此外，对总体性原则的提升，很可能成为形而上学渗透的支点。事实上，斯大林和恩格斯一样，都强调了相关性观念，而不是分裂观念的重要性。

此外，原则1（普遍相互依赖）和原则2（上升力量的第一性）之间没有建立起真正的内在一致性。相互作用如何反过来被新事物的载体力量的行动所决定？在斯大林那里，结构问题（过程各环节的结合）与趋势问题（矛盾主要环节的决定性作用）仍然是分离的。这反映出，共产主义的未来并没有在指导无产阶级专政的当前矛盾运动中被清晰地构想出来。在斯大林的强有力思想中，没有抓住两条道路斗争的方法。从这个意义上说，斯大林的表述重复了我们发现恩格斯所做的对矛盾原则的削弱。

\subsection{列宁}

列宁在关于黑格尔《逻辑学》的笔记中，在评论关于绝对理念的一章时，开始列举他称之为“辩证法的要素”。

显然，这个清单旨在对黑格尔一个相当复杂的章节进行分析性澄清。它不代表一个连贯的理论体系。

列宁列举了十六个要素，但他自己指出，它们并非都具有相同的地位。他含蓄地批评恩格斯，例如，他指出第16点（量变到质变）只是第9点（每个规定性向相反的规定性的转化）的一个例子。其他一些点实质上是相同的。

例如，第14点（表面上回到旧事物）的内容实际上已经在第13点（在更高阶段重复较低阶段的某些特征）中出现。（列宁将这种表面上回到旧事物称为否定之否定。在我们看来，这个称谓是模棱两可的：表面上回到旧事物确实是历史发展的一条规律，中国人称之为螺旋式发展规律。我们将在其他地方解释其含义。让我们立刻说，这与否定之否定毫无关系。）

最终，我们可以将列宁的十六个要素归纳如下：

\begin{enumerate}
    \item 关于认识过程的唯物主义论点：
    \begin{itemize}
        \item “考察的客观性”（第1点）；
        \item 认识过程的无限接近性（第10、11、12点）。
    \end{itemize}
\end{enumerate}

我们将注意到，这两个特征分别指向了反映的隐喻和渐近线的隐喻。唯物主义认识既是现实运动的反映，也是对这种运动的趋向性逼近。它既复制了现实运动，又趋向于与它等同，但永远无法完全达到。认识是一个运动中的影像。

作为物质过程，它根据其精确性（反映，认识的绝对性）和其不精确性（趋向性，认识的相对性）而分裂。这种分裂的内在规律是以渐近的方式消解于统一中。但二者永远不会融合为一，运动中的影像始终与其客体分裂。

列宁在这里重申了《唯物主义和经验批判主义》的主要论点。（“绝对真理是由发展中的相对真理的总和所构成；相对真理是独立于人类的客体的相对正确的反映；这些反映变得越来越正确；每个科学真理，尽管其具有相对性，都包含绝对真理的要素，——所有这些对于任何思考过恩格斯《反杜林论》的人来说显而易见的命题，对于“现代”认识论来说都是天书。”\footnote{列宁，《唯物主义和经验批判主义》，莫斯科出版社，1952年，第359页。}）

正是在这一点上，唯物主义和辩证法在认识论中相互关联：认识是反映-过程（参见《唯物主义与唯心主义》）。

\begin{enumerate}[resume]
    \item 普遍相互依赖的论点：
    \begin{itemize}
        \item “事物之间多样的和普遍的联系”（第2点）；
        \item “每个事物都与每个其他事物相联系”（第8点）。
    \end{itemize}
    \item 普遍运动的论点（第3点）。
    \item 作为现象内在本质的矛盾论点：
    \begin{itemize}
        \item “事物内部的矛盾倾向”（第4点）；
        \item “事物作为……对立面的统一”，以及“这些对立面的斗争”；
        \item “每个规定性、性质、特征、方面、特性向另一个规定性（向其对立面？）的转化”（第9点）。
    \end{itemize}
    \item 螺旋式发展的论点（第13、14点）。
\end{enumerate}

其余的论点实际上，正如列宁自己所说，只是关于矛盾论点的例子，分别应用于分析和综合（第7点），形式和内容（第15点），量和质（第16点）。如果我们抛开对纯粹唯物主义论点的重申（第1、10、11、12点），我们最终得到了辩证法的四个原则：

1. 相互依赖（总体性原则）。
2. 任何现实都是运动、过程。
3. 矛盾是过程的本质。
4. 螺旋式发展。\footnote{列宁，《哲学笔记》，社会出版社，第209页及以后。}

这四个原则既不是恩格斯的，也不是斯大林的。它们显然围绕着列宁所认为是黑格尔的决定性贡献而分组：将现实构想为矛盾的运动。此外，在这次分析性的、没有真正排序的列举之后，列宁提出了一个唯一的原则，并用框线突出其重要性：

\begin{quote}
“可以简要地把辩证法定义为对立面的统一理论。这样，就能抓住辩证法的核心，但这需要解释和发展。”\footnote{列宁，《哲学笔记》，社会出版社。}
\end{quote}

在1915年的《论辩证法问题》片段中，列宁有力地重申了这一观点，并这次不再从对立面统一的角度，而是更根本地从统一分裂的角度来阐述：

\begin{quote}
“一分为二，以及对它的矛盾部分……的认识，是辩证法的实质（基本特征之一、基本特点之一，甚至是根本特点）。”\footnote{同上。}
\end{quote}

在世界大战的背景下，列宁对神圣同盟、帝国主义领导下的阶级合作、第二国际的破产感到愤怒，他将辩证法作为一种关于分裂和对抗的思想来掌握。他首次提出了革命辩证法的真正母体，即“一分为二”原则。

这确实为工人运动的分裂、为对第二国际的必要摧毁性批判做了哲学上的准备。正是两条道路斗争在世界大战考验下的必然爆发，成为辩证法思想决定性前进的历史基础。

同样，中国共产党与苏联修正主义的斗争，也反映在中国国内激烈的哲学论战中，将列宁的“一分为二”原则与那种没有原则的“和平共处”和投降主义修正主义论点对立起来，其哲学核心是声称辩证法的主要格言是“合二而一”的“规律”。

我们在这里清楚地看到，只有历史的狂暴才能在认识运动中创造性地发挥作用。列宁不仅超越和批判了恩格斯局限的表述，不仅将矛盾置于对现实的一切理解的中心，而且毫不犹豫地将矛盾对同一性的第一性推向最极致：

\begin{quote}
“对立面的统一（同一，同一性，等同）是有条件的、暂时的、相对的。相互排斥的对立面之间的斗争是绝对的，如同发展和运动是绝对的一样。”\footnote{同上。}
\end{quote}

斗争是辩证法思想唯一的绝对原则：这就是作为反抗哲学的辩证法的本质。

列宁留给其继承者的问题是：如何从这种反抗出发，从集中了其所有方面的唯一论断，即对立面统一的规律出发，来阐明辩证法的原则？

完成这项任务，就是提供列宁所说的“解释”和“发展”，使得人们能够抓住包含在对立面统一这一个原则中的“辩证法的核心”。牢牢抓住对立面统一，即“一分为二”原则，作为矛盾论的公理化核心，并从此出发发展出从属论点的链条，这就是列宁的哲学指令。

\section{论矛盾}

在《矛盾论》的第一句话中，毛泽东就重复了列宁关于辩证唯物主义有一个唯一基本原则的观点：

\begin{quote}
“事物发展的根本法则，即对立统一的法则，是唯物辩证法的根本法则。”
\end{quote}

毛泽东明确地将他的论著中讨论的问题，作为对一个“广泛问题圈”的发展，而这个“问题圈”完全从属于对唯一基本原则“对立统一”的研究和阐明。毛泽东所遵循的阐述方法是从一般到特殊，即从矛盾规律的普遍性，到分析一个特定矛盾运动的范畴（矛盾的主要方面，对抗性在矛盾中的位置）。

从我们感兴趣的角度来看，也就是从辩证论点的有序重构角度来看，我们可以认为毛泽东将对立统一的唯一原则注入了五个本质性的辩证论点中。

这五个论点是：

\begin{enumerate}
    \item 任何现实都是过程。
    \item 任何过程，归根结底，都可以归结为一个矛盾体系。
    \item 在一个过程（即一个矛盾体系）中，总是有一个主要矛盾。
    \item 任何矛盾都是不对称的：换句话说，矛盾的一个方面总是在整个矛盾运动中占主导地位。这就是矛盾主要方面的理论。
    \item 存在不同类型的矛盾，其解决方式也不同。在这方面，主要区别在于对抗性矛盾和非对抗性矛盾。
\end{enumerate}

因此，毛泽东完成了列宁制定的计划：发展“辩证法的核心”，即“一分为二”这一公理。这一发展是对反抗教条主义和反抗修正主义的马克思列宁主义的哲学综合。

中国革命的历史条件要求对马克思列宁主义做出创造性贡献。1925-1927年城市起义的失败，以及随后被迫撤离井冈山根据地和长征，表明十月革命的理论教训不足以在中国就夺取政权问题制定和实施一条可行的路线。

此外，还必须在党内摧毁教条主义的拥护者，特别是在革命战争过程方面。

军事问题确实在深化哲学问题方面发挥了决定性作用。我们知道，斯大林通过推广列宁所总结的城市起义的特定规律，支持了进攻的军事理论，并鄙视克劳塞维茨提出的防御具有战略优势的论点。

同样，中国共产党的一个流派在1925-1930年间也支持对城市的进攻路线，或是在没有任何退却的情况下在原地防御解放区。

诚然，全面进攻、不可能停顿的规律适用于起义：巴黎公社从反面历史地证明了这一点，十月革命则通过胜利证明了这一点。但起义的规律涉及的是一种通过力量的极速积累而实现的瞬间性突破现象。它们与战争的过程相去甚远，后者辩证地涵盖了对时间和空间的巨大考量。

毛泽东对战略防御进行了大量论述，这很有特点，他运用了严谨的辩证法，来处理对抗力量在时间上的部署。恰恰是《中国革命战争的战略问题》这本小册子的一个目标，就是反对对苏联军事理论的盲目模仿：

\begin{quote}
“还有一种观点也是错误的，我们早已驳斥过它。他们说……只要照搬苏联国内战争的经验，照搬苏联军事机关颁布的军事条令就行了。他们不了解，这些经验和条令反映了苏联国内战争的特殊性，如果生搬硬套，不加任何修改，就又会犯‘削足适履’的错误，再一次招致失败。”\footnote{毛泽东，《中国革命战争的战略问题》，《毛泽东著作选集》第一卷，人民出版社，1991年，第195页。}
\end{quote}

战争与哲学之间自古以来就有特殊的联系。在关于军事问题的思想斗争背后，我们很可能发现毛泽东对斯大林提出的不满：他是一个糟糕的辩证法家，而且对德国哲学的研究太少。毛泽东和中国共产党只有付出进行无情思想斗争的代价（反对陈独秀、李立三、王明等），才能锻造出一条新路线的积极要素。

党本身的分裂将两条道路斗争的概念提升到理解政治现象的首要原则。因此，辩证系统化变得极其重要。

在社会主义建设时期，这方面的问题更是得到了加强：关键点是对苏联资本主义复辟的总结。这一总结证明，斯大林所选择的道路不能作为榜样，因为它未能阻止赫鲁晓夫的资产阶级政变。同样，在这里，也必须与守旧者进行斗争，并将路线问题的核心，不是放在生产力的发展上，而是放在无产阶级专政下的阶级斗争上。不是放在经济基础的积极统一上，而是放在上层建筑，以及党本身，分裂为无产阶级道路和资产阶级道路上。

从井冈山到文化大革命，毛泽东思想以逆流、以分裂的工作而形成。因此，它以越来越精准的方式展开了对立面统一原则的后果。它是一种卓越的反叛思想，是一种反抗的反抗思想：辩证思想。

这五个原则相互关联，仿佛要深入到对抗的中心，它们很好地代表了革命认识，即朝向行动的运动，是一种打破现有统一的立场，以便根据分裂所产生的能量，胜利地掌握未来的统一。

让我们沿着这条路走下去。

\subsection{第一个原则：任何现实都是过程}

这个原则表面上非常简单。它将现实描述为运动，它断言事物的内在本质，除了其自身转化的规律之外，什么也不是。这个原则属于辩证法的赫拉克利特传统，即：“一切都在变化。”在这里，反对《传道书》中那种沉闷的服从智慧的反抗在思想中活跃起来。它用永远崭新、高高升起的红色太阳来对抗“太阳底下没有新鲜事”，在这面旗帜下，反叛的劳动者无限的肯定性希望产生了突破。

然而，这个简单而暴力的原则是一个受到威胁的原则，必须在激烈的阶级斗争中不断重新夺回，因为它划清了与对抗性倾向——形而上学倾向——的主要分界线。这个原则的本质在于断言，现实的任何给定状态在原则上都是暂时的，换句话说，事物的规律永远不是平衡，也不是结构，而是所有平衡的打破，因此是现有事物状态的必然发展和摧毁。

这是这个第一原则的纯粹革命意义：它站在一个永远不可能是保守主义的立场。从某种意义上说，它摒弃了一切客观主义：在一个时刻呈现出的现实，其本质只是这种现实正在消解并转化为另一种现实的运动。严格来说，现实并不属于客体范畴。客体是作为状态或形态而为人所知的。然而，任何状态或形态的内容都是其自身不间断的变形过程。

从这个意义上说，客体与过程对立，就像形而上学与辩证法对立一样。形而上学，归根结底是同一性理论，它被一种强大的保守运动所驱动。它是一种对现实某种状态的守护事业。它是一种防御性的阶级立场，它紧紧抓住表象，以掩盖其所守护的现实不断逃离它的本质运动。

马克思从一开始就断言，这就是辩证法的直接政治使命：站在摧毁一边，反对形而上学的守护：“辩证法在其合理的方面，是统治阶级及其教条主义思想家们所不能容忍的，因为在对现存事物的肯定性理解中，它同时包含对现存事物的必然否定、必然灭亡的理解；因为它把任何已经形成的形态都看作是暂时的，看作是运动中的一个瞬间，所以它不为任何事物所惑；因为它本质上是批判的和革命的。”\footnote{马克思，《<资本论>第一卷德文第二版跋》，《马克思恩格斯全集》第23卷，人民出版社，1972年，第24页。}

辩证法第一原则的批判意义必须反过来受到保护，免受反动形而上学的压力。哲学修正主义恰恰在于表面上承认任何现实都是过程，但又固化了一个过程概念，这使得制约其转化的规律变成了新型的形而上学不变性。

在黑格尔那里，列宁所赞扬的对“现实是过程”这一原则的承认，在主导的唯心主义重压下也分裂为二。事实上，对于黑格尔来说，过程总是对一个简单术语的发展，结果是过渡形态并非真正被摧毁，而是在一种循环运动中被来回穿梭，这种运动只是为了在更高层次的同一性中保留它们。

黑格尔试图从一个不可分解的初始姿态中产生运动的必然性。结果是形而上学从内部侵蚀着辩证法，从而在他的著作中产生了断裂和不一致，这构成了两种倾向斗争的真正症候。（我们将在关于黑格尔辩证法的论述中指出，对黑格尔的马克思主义解读本质上是一种分裂的解读：它在辩证过程的断裂和不一致中，发现了唯物主义和唯心主义倾向的冲突。

恩格斯也说了这一点：“尽管黑格尔以如此的正确性和天才看待了许多特定的关系，但由于上述原因，细节也经常变成拼凑、人为、构建，简而言之，是对真理的歪曲。”\footnote{恩格斯，《空想社会主义和科学社会主义》，作品全集，第三卷，第136页。}）

这是因为简单开端的观念是一种典型的形而上学，即保守的预设。在马克思列宁主义毛泽东思想中，实践第一性的痕迹是，理论本身永远不会开始任何事情。

不可能有理论上的零度：认识的运动总是早已开始，因为它包含了实践作为其一个环节。而作为直接现实的实践，不断地预设自身。正如列宁所很好地指出的：

\begin{quote}
“在自然界和生活中，存在着‘走向虚无’的运动。但‘来自虚无’的运动，无疑是没有的。总是从某个东西开始的。”\footnote{列宁，《哲学笔记》，社会出版社，第127页。}
\end{quote}

不仅一切都是过程，而且 a) 所有过程都与其他过程相联系，b) 对一个过程的任何认识，作为其发展的一个时刻，总是早已开始。

否认这一点，就是：

\begin{itemize}
    \item 或者像黑格尔那样，是彻头彻尾的唯心主义者：认识开始了，因为理念产生了作为整体的现实；
    \item 或者假设存在一个绝对的立足点，一个对现实过程的均衡暂停，在那里理论找到了一个不动的客观起点，因此可以脱离正在进行的过程的“已然存在”。因此，理论不过是这种假定平衡的守护者：这是形而上学的保守功能。
\end{itemize}

如果我们现在来看阿尔都塞和 D. 勒库尔坚持将运动的辩证原则表述为“没有主体也没有目的的过程”，我们会发现这包含了一种限制性的诡辩，以及一个阵营的选择。这种全面的客观主义掩盖了所有变化都是从构成其矛盾本质的对抗性力量的角度来运作的。这等于是将运动的论点与另一个论点分开，即矛盾的一个方面，而不是两个方面，是承载过程的正确性，即其未来。

阿尔都塞在这方面进行了一场令人不安的交易，混合了刻意的沉默和不恰当的替换，涉及列宁的《哲学笔记》。他首先做出了一个正确的观察，即列宁摒弃了绝对起源、理念的自我运动第一性等所有问题。我们完全同意，思辨性的“开端”概念——黑格尔自己也对此进行了片面的和不彻底的批判——是用来将辩证法缝合到唯心主义上的。这种联系必须被打破，而黑格尔确实在打破它（在局部的辩证序列中）和维持它（在全局的体系中）之间摇摆不定。

根据阿尔都塞的说法，列宁在《逻辑学》中关于绝对理念的章节中，所着迷的是，在起源的废墟中，将绝对等同于辩证方法，即等同于过程。这并非不正确：在黑格尔那里，确实有可以为我们的第一原则（任何现实都是过程）提供养料的东西。

但这里出现了偏差：对阿尔都塞来说，过程（即作为运动的现实）是唯一的绝对，这意味着列宁在黑格尔那里找到了“必须……消除所有起源和所有主体，并说：绝对的就是没有主体的主义”的“确证”。\footnote{路易·阿尔都塞，《列宁面对黑格尔》，马斯佩罗出版社，第88页。}

阿尔都塞的“没有主体的主义”范畴在这里有什么用呢？列宁在他对绝对理念的评论中没有提及主体问题。甚至更好的是：当他——在其他地方——讨论黑格尔专门讨论主体范畴的章节时，列宁也没有以任何方式“消除所有主体”。他感兴趣的是抓住主观与客观的辩证法，即矛盾运动。对于列宁来说，问题是要依靠黑格尔来终结主体和客体范畴的片面性，只要它们被分开（形而上学操作）或其中一个被消除（绝对唯心主义或机械唯物主义）。

\begin{quote}
“恩格斯《自然辩证法》中非常好的一段话，其中‘认识’（‘理论的’）和‘意志’（‘实践的’）被描述为消除‘客观性’和‘主观性’的片面性的两个方面、两种方法、两种手段。”\footnote{列宁，《哲学笔记》，社会出版社，第198页。}
\end{quote}

因此，列宁的问题不可能成为“消除”主体，这是一种回归到客观性片面性的倒退。问题是要同时反思这两个范畴（主体和客体）在一个过程的整体运动中的分裂和相互作用，而不排除主观因素可能成为这个运动的关键。列宁以最接近黑格尔的方式阐述了这个计划：

\begin{quote}
“如果我们从逻辑上来看主体与客体的关系，我们也必须考虑到具体主体（=人的生活）在客观情况下的普遍存在前提。”\footnote{列宁，《哲学笔记》，社会出版社，第192页。}
\end{quote}

阿尔都塞的整个操作归结为提出一个等式：起源=上帝=主体，并得意洋洋地得出结论，通过消除起源和上帝，列宁在黑格尔之后构建了没有主体的过程概念。

但阿尔都塞在这里只是慷慨地将他自己的辩证法短视归于列宁。起源（奠基者、先验者……）不是主体，而是其唯心主义谓语。消除起源或上帝——这是基本的唯物主义者所要求的——的结果仅仅是主体的概念分裂了，其唯心主义表述被强行与来自宗教或法律的意识形态概念联系在一起，而其唯物主义表述则在过程的范畴中得到体现。是一个有主体的过程。

没有什么比阿尔都塞顽固地将列宁-黑格尔关系归结为消极内容更令人震惊的了，归结为两个根本上配对的消极内容：对康德的批判和对主体范畴的批判。（D. 勒库尔作为一名好学生，强调并加剧了这种顽固：“他[列宁]在他的笔记中保留的，本质上——如果不是排他性地——是黑格尔批判康德哲学的部分。”\footnote{多米尼克·勒库尔，《一个说谎者及其赌注》，马斯佩罗出版社，第51页。} 无法更强烈地表达如此强烈不正确的事情了。）

然而，如果坚持只看关于绝对理念的笔记，有两个坚定的积极思想让列宁着迷，而阿尔都塞对这两个思想只字未提——这确实是可耻的。这两个思想是：

\begin{enumerate}
    \item 理论与实践的辩证统一，实践被包含在认识的运动中，“实践标准”是真理的唯一标准。
    \item 对立面的统一是“辩证法的核心”。
\end{enumerate}

这都是两个辩证法思想。将它们简化为唯物主义显而易见的消极平庸性：没有起源也没有上帝，是一种残缺。将它们视为排斥主体，是一种伪造：因为正如我们在列宁自己的文本中看到的，这两个思想的关联恰恰被注入了客观和主观的矛盾过程中。

革命认识——列宁在黑格尔身上看到其“萌芽”——是同时将阶级实践、反抗纳入其运动，并回归其中以组织其自觉胜利的东西。它是一个居于阶级斗争的客观过程和以计划为导向的无产阶级革命的主观实践之间的中心化的辩证中介。

黑格尔的“主观”术语仍然适用于描述这一进程：从自在到自为。这些自为和自在的范畴，在阿尔都塞的“没有主体的主义”中是无法理解的，而列宁则赞扬它们“极端的正确性和精确性”。\footnote{列宁，《哲学笔记》，社会出版社，第193页。}

直到今天，恩维尔·霍查仍然没有找到更好的概念来处理工人实践的客观条件与其革命目标之间的关系问题：

\begin{quote}
“历史已经证明，没有自己的党，工人阶级，无论其生活和行动处于何种条件，都无法自己获得阶级意识。将工人阶级从‘自在阶级’转变为‘自为阶级’的是党。”\footnote{恩维尔·霍查，《在阿尔巴尼亚劳动党第六次代表大会上的报告》，《奈姆·弗拉什里》出版社，第234页。}
\end{quote}

在人类解放的普遍过程中，被剥削阶级始终是历史的主体，因为创造历史的是群众。如果没有上帝，没有起源，那么确实存在创造性的主体：

\begin{quote}
“人民，只有人民，才是创造世界历史的动力。”\footnote{毛泽东，《论联合政府》，《毛泽东著作选集》第三卷，人民出版社，1991年，第1031页。}
\end{quote}

但奴隶和农奴仍然是客观主体，是部分未分裂的主体（最多只是胚胎性的、空想性的），分为自在主体和自为主体。无产阶级通过马克思主义与现实工人运动的融合，即作为客体-主体的阶级与作为主体的阶级的分裂关联，来实现对其主体角色的主观占有，这种关联体现在这个前所未有的过程中：阶级党。

从《宣言》开始，很清楚，对于马克思和恩格斯来说，共产主义者不能将他们的认同建立在一个“没有主体”的革命过程中：如果仅限于现实运动，他们只是所有国家工人党“最先进的部分”。他们正是通过主观因素与工人运动进行辩证，因为“他们比其余无产阶级群众优越的地方，就在于他们了解无产阶级运动的条件、进程和一般结果。”\footnote{马克思，恩格斯，《共产党宣言》，《马克思恩格斯全集》第一卷，人民出版社，1995年，第285页。}

D. 勒库尔说，这是一个没有主体也没有目的的过程。那么共产主义者是什么？事实上，共产主义者不过是通过掌握历史规律和目标，使得“工人党”（战斗阶级）成为历史过程主体的那个东西。

但我们不会惊讶地看到“理论实践”的发明者，对列宁因黑格尔发现实践是认识的一个矛盾环节而感到的热情视而不见。至于主体的消失，它非常适合那些回避法国无产阶级革命根本问题的人：无产阶级的主观状态问题。

更具体地说，当有人向他提出唯一值得回答的明确问题时，他回避、退缩，像政治乌贼一样，将自己淹没在墨水中：阿尔都塞，你作为其中一员的法共，今天在法国是“工人阶级的党”吗？它是作为其革命历史主体的无产阶级吗？如果不是，那么你那关于“没有主体的主义”的被动理论，到底能为反抗篡位者的哪一次反抗提供正确性呢？

对于阿尔都塞来说，造反反动派从未真正有理，特别是反抗法共中的反动派。以客观主义为掩护，否认无产阶级能够成为历史主体，就是为那个声称代表它的资产阶级客体——修正主义党——留下空间。

以反对唯心主义目的论为借口，否认无产阶级承载着共产主义计划，即人类史前史的终结，就是支持社会主义建设中的资产阶级倒退力量，支持那些正是有意让无产阶级专政脱离对其终结——无阶级社会，国家的消亡——的任何清晰认识的人。

从根本上说，“没有主体也没有目的的过程”理论，通过其所固定的客观性，取消了第一原则（运动原则）。它使得将演变作为客观结构的历时性结果，优先于将运动作为斗争和新事物。它站在现有平衡的保守一边，任何反抗群众作为历史主体的突然涌现都不得来扰乱它。它本身就是形而上学渗透到辩证话语中的典型例子：这正是哲学中的修正主义的定义。

因此，辩证法的第一原则必须顽强地得到捍卫，以对抗其修正主义的腐败（将过程视为客观的结构同一性）和其公开的形而上学腐败（将过程视为复杂事物从简单事物的自我产生，假设存在一个起源）。

坚守这一原则，决定性地关系到哲学中的阶级立场。

\subsection{第二个原则：任何过程都是一个矛盾的集合}

这里划定的界线，使得辩证法与进化论彻底对立。运动不是通过连续环节的线性产生而运作的：过程中的任何一个环节本身都处于分裂状态，而运动的动力正是这种分裂。对形而上学同一性原则的批判必须延伸到这样的论点：现实的一个过渡状态，其存在恰恰是过渡本身，也就是它所发展的内部裂变。

仅仅说事物处于运动中是不够的，还必须认识到，“事物”这一概念本身，不是遵循同一性逻辑，而是遵循分裂逻辑。现实不仅不会消解在某个状态或平衡的单一性中，而且单一性本身也只能被理解为分裂。运动不是一连串的统一体，而是一团分裂的交织。

借用列宁的话，我们将说：只有将对现实是过程的认识，延伸到对“一分为二”原则的认识，才是辩证法家。

“一分为二”不是从“一”中产生“二”的原则。“一分为二”意味着：只有分裂的同一性。现实不仅是过程，而且过程是分裂。现实不是聚合，而是分离。即将发生的事情是分离开来的。

一个过程越是发展，其各环节之间冲突性的增长就越是得到肯定，分裂原则就越是倾向于将最初矛盾的每个环节分离开来，特别是上升的环节，那个“胜利的”环节：因为正是它将承载新的统一，即新的分裂。

例如，无产阶级党，是其阶级存在的集中形式，即其与资产阶级对抗的集中形式，也是其胜利（建立无产阶级专政）的决定性工具，它同时也是一系列不间断分裂的场所，是分裂的首要场所。

\begin{quote}
“老黑格尔已经说过：‘一个党只有通过分裂自身并能够承受分裂，才能证明自己是一个胜利的党。’无产阶级运动必然要经历不同的发展阶段；在每个阶段，一部分人会停下来，不再继续前进。这正是为什么‘无产阶级的团结’到处都以不同党派群体的形式实现，这些党派群体进行着你死我活的斗争，就像罗马帝国在最严酷的迫害下，基督教派系所做的那样。”\footnote{恩格斯，《致倍倍尔信》，1873年6月，作品全集，第二卷，第453页。}
\end{quote}

（顺便再次赞扬一下黑格尔辩证直觉的深刻。这句令人钦佩的引文来自《精神现象学》。黑格尔在那里，在讨论启蒙运动时，发展了夺取政权的逻辑（这里是意识形态上对宗教信仰的完全胜利）以及新两条道路斗争立即分裂胜利阵营的逻辑。\footnote{黑格尔，《精神现象学》，奥比尔出版社，第二卷，第123页。}）

对于毛主义者来说，记住恩格斯对抗那些哭喊着不惜一切代价要统一的人，对抗那些坐在高凳上对“小团体争吵”发布超脱判决的小资产阶级，肯定分裂的创造性价值，并坚持通过这些激烈的争吵，正在锻造一种新型的“工人团结”，即无产阶级革命的一个新进展，这是一种强大的鼓舞。

今天在法国，各种团体的冲突多样性反映了运动的真实状态，它们的思想冲突的激烈性，远非一种令人沮丧的人为现象，而是未来将要出现的场所。浸入这些冲突并明确地站队：其余的都不过是虚假的万众一心、无效的民粹主义或修正主义的沉睡。

一分为二。

认识并坚持这一格言并不容易。在中国，“一分为二”原则曾是激烈论战的焦点，这些论战直接关系到文化大革命的思想准备。正是在1964年，杨献珍提出了“合二而一”作为辩证法的第二个核心。

杨献珍的论点是，他将对立面统一的普遍原则分裂，分别强调了“统一”（合二而一）和“对立”（一分为二）；然后，最重要的是，他让这两个论断等同。我们在这里立即再次发现了所有修正辩证法的形而上学操作的动力：即平衡的假设（在这里，是两个矛盾论点之间的平衡），通过这种偏见，同一性再次压倒了矛盾。（关于所有这些，参见1971年5月24日《北京周报》的文章：“‘合二而一’是复辟资本主义的反动哲学”。）

为了证明统一与对立之间的平衡是正确的，杨献珍必须恰恰终结矛盾的普遍性，也就是终结任何事物不可避免地分解为矛盾部分的过程。为此，他：

a) 将（对立面的）统一赋予了一个可分离的、与两个方面都不同的积极内容：他谈到对立面之间的“共同点”，统一是这些“共同点”的存在和总和的结果。因此，斗争的动态性被两个兼容的总体性之间的对抗所取代：统一变成了两个集合的完整交集。

辩证法论点对此反驳道，对立面的统一不过是它们的关联、互补性，也就是它们的斗争、相互排斥。当然，没有资产阶级就没有无产阶级，没有马克思列宁主义者就没有修正主义者。但资产阶级和无产阶级之间，修正主义者和马克思列宁主义者之间，除了通过他们不可调和的冲突性对立来定义一个历史过程——阶级斗争、工人运动中两条道路的斗争——的分裂统一之外，没有任何共同之处。对立面的统一是一种关系，而不是一种同一性。

对立面的交集是空的。它只有在我们将两个方面视为在与第三方的对立中成为一个整体时才会“填满”，但我们已经改变了过程：无产阶级和民族资产阶级，作为对立面，没有任何共同之处。然而，人民，在他们与帝国主义侵略者的对立中，可以同时包含无产阶级和民族资产阶级。

但此时所涉及的过程不再是最初的过程（无产阶级革命），而是一个新的过程（民族解放战争），其内容是一个新的矛盾（人民/帝国主义侵略者）。然而，从这个过程的角度来看，“二”并不比之前更“合一”：人民和帝国主义侵略者之间没有任何共同之处。

b) 杨献珍认为存在不可分割的环节，无法解开的联系。

辩证法相反地肯定了分裂原则的普遍性：

\begin{quote}
“在人类社会中，如同在自然界中一样，一个整体总是分裂为部分，只有其内容和形式随着具体条件而变化。”\footnote{毛泽东，《在中国共产党全国宣传工作会议上的讲话》，《毛泽东著作选集》第五卷，人民出版社，1977年，第420页。}
\end{quote}

c) 杨献珍认为综合（合二而一）必须强制性地补充分析（一分为二）。

辩证法不接受这种类型的综合，它是一种典型的机会主义概念。一个新过程的出现（一个新的对立面统一）是通过前一个矛盾的环节的消失而实现的，也就是说，是在一个方面摧毁另一个方面之后，这种摧毁必然导致胜利方面自身的分裂。

正是这种分裂将定义和支配新的过程。综合是摧毁/分裂的过程。资产阶级摧毁了封建秩序，但它所领导的资本主义秩序立即分裂为资产阶级国家/无产阶级革命的矛盾。

辩证法的综合概念，就是产生一个新的分裂，仅此而已。

因此，杨献珍的哲学进攻确实旨在削弱第二个原则：矛盾原则，通过施加一个平衡的界限，即一个等同的原则（“合二而一”），这个原则反过来又被转化为内容的同一性、不可分割性和综合。

其背景是：国际共产主义运动的分裂，以及与现代修正主义及其堡垒——苏联国家——进行你死我活斗争的广度和影响。

杨献珍是调和对立面的哲学家，其掩护是成为对立面统一的哲学家。

为了衡量在中国人所称的“哲学战线”上进行斗争所必需的警惕性，让我们看看这个事实：在六十年代，杨献珍在他的“合二而一”概念是在党的高级党校里发展的。

\subsection{第三个原则：在任何过程中，都有一个主要矛盾}

在我们目前的讨论阶段，人们可能会认为，作为分裂交织的现实不包含任何特殊领域，而只有这种交织的一般体系。从某种意义上说，这是对的：所有矛盾都是相互依赖的。我们已经看到恩格斯和斯大林都充分发展了这一点。

然而，从另一种意义上说，我们可以理解具有自身辩证一致性的现实领域。它们是矛盾体系，其质的规定性是由其从属于一个主要矛盾所决定的。只有当我们设想一个矛盾体系，其相互依赖性通过其对体系中一个主要矛盾的质的从属关系而得到调控时，我们才能有意义地谈论一个过程。

例如，一个资本主义社会的历史运动可以被认为是一个特定的质的规定过程，只要我们在其中找到资产阶级/无产阶级矛盾作为主要矛盾。一场特定的工厂斗争是由工人/资本家这个主要矛盾所规定的。很清楚，主要矛盾决定了过程的质的性质，其一般的统一类型，但过程本身也包含许多其他矛盾。

例如，一次工人罢工必然会包含，并常常加剧，工人之间的许多矛盾（罢工者/工贼，革命工人/工会主义者等）。这些除了主要矛盾之外的其他矛盾通常被称为次要矛盾。但次要矛盾发展的内在衡量标准，取决于它们在所设想的过程的观点上，与主要矛盾的衔接。

例如，罢工中工人之间的矛盾，只要它们影响主要矛盾（工人/资本家）的发展（削弱工人阵营，表现出极左倾向，工会清算等），就属于“罢工”过程。在运动过程中，主要矛盾和次要矛盾之间的辩证联系不断变化，并以相互作用的方式变化。

例如，很明显，资本家会竭尽全力加剧工人之间的某些矛盾，因为反过来，这种加剧会改变主要矛盾中各方面的关系，使之有利于他。但毛主义者也可以攻击某种“软弱”的工会统一：工人之间严重的矛盾可以成为加强他们阵营的因素，例如，如果它们以胜利的方式使一贯的革命方向与修正主义方向对立起来。

突然，过程的质的性质被改变了；它将不再主要被定义为“罢工”，而是作为两条道路的斗争。它的空间将不再是经济的，而是政治的。罢工所属的整体过程将不再是围绕剩余价值率的阶级冲突，而是反修正主义的新革命组织的建立，等等。

同样，对主要矛盾原则的静态认识可能属于形而上学。仅仅承认任何过程都由一个主要矛盾定性是不够的。还必须认识到，作为整体的矛盾体系也是运动的，是过程，而对现象的理解则在于主要矛盾和次要矛盾之间不断变化的相互关系。

在这方面，源自68年思潮的“群众主义”意识形态擅长于削弱辩证分析。在《解放报》或任何一份“极左”报纸上，到处都会谈论“正在斗争的工人”、“运动”、“雇员与老板的斗争”、“反对体制的生态斗争”等等。

这是对同一种“主要”矛盾的无限重复：群众对抗权力，却没有看到，把异质和深度分裂的整体塞进同一个袋子（“运动”）里，会妨碍我们理解新事物在哪里，以及在什么条件下，作为某个过程的次要矛盾的主要发展而出现。

然而，新事物总是以这种方式出现的：它从来不是“群众”，也不是“运动”整体地孕育它，而是它们之中与旧事物分裂出来的东西。例如：工人反抗中摆脱工会主义的东西，或葡萄牙人民运动中反对主要资产阶级力量（法西斯主义和社会法西斯主义）控制的东西。

这种对主要矛盾的冻结，表面上是“极左”的（总是同样的狂热群众对抗同样的权力、一成不变的体制），与修正主义完全趋同，后者也试图通过鼓吹人民“团结”对抗现有权力来保护其政治立场。

哪个毛主义者没有被一个法国总工会的走狗称作“分裂分子”？是的！我们支持“一分为二”。我们支持新事物在分裂中成长。我们既不想要被神化和模糊的、无效和重复的“极左”群众，也不想要修正主义的统一，那是一个阴险独裁的门面。今天，无产阶级的东西首先是分裂、斗争，并使得“运动”内部微小的裂缝成长为主导。

只有将对主要矛盾原则的认识推向对次要矛盾成为主要矛盾的认识，才是辩证法家。

\subsection{任何矛盾都是不对称的：它有一个主要方面}

正如其主要矛盾决定了一个过程的质的性质一样，每个矛盾也由一个主要方面来定性。这里的“方面”一词不应产生幻觉：它不代表一种形而上学类型的不可分解的统一体的重新出现。“方面”除了在其所参与的分裂过程中之外，没有任何现实性。

资产阶级的辩证法定义是其与无产阶级对立的所有实践（剥削、压迫、反革命）的集合。一个阶级永远不会先于阶级斗争而存在。存在就是对立。一个方面的存在，完全是在其与另一个方面的矛盾关系中给定的。说其中一个方面是主要的，就是说正是从它在分裂过程中的主导地位，我们才能确定分裂本身整体的质的性质。

例如，法国是一个资本主义国家，这不仅因为主要矛盾是资产阶级/无产阶级矛盾，而且因为在这个主要矛盾中，资产阶级是矛盾的主要方面：资产阶级占主导地位，也就是说，它仍然能够以对其有利的方式，确定其与无产阶级对抗的方式和框架。

很清楚，关于主要方面的理论，就像关于主要矛盾的理论一样，其内容是主要方面与非主要方面之间关系的运动。辩证法所研究的，与其说是对主要方面的识别，不如说是次要方面如何成为主要方面，以及其辩证的关联：主要方面如何成为次要方面。

但是，这恰恰产生了一个问题。在这种情况下，如果“主要”的本质是转化为其对立面（次要），那么“主要”这个词指的是什么？更重要的是：如果主要和次要相互转化，我们是否应该理解辩证法的最高规律只是一种简单的置换原则，一种简单的位置交换？

这个问题具有重大的哲学意义。置换、交换规律，确实是所有结构主义意识形态的本质动力。而这些意识形态的政治影响是众所周知的：如果现实的运动消解于一种组合分析中，那么可以确定，这种运动的本质在于其不变性，即制约置换的规则。

从这个意义上说，任何新事物都很大程度上是表象：各方面的位置移动，却使得潜在的交换结构保持不变。表象的流动性指向一个封闭的体系。所有结构性思想的本质保守主义，在这里有将辩证法变成其对立面——形而上学——的风险。

今天，这种观念最活跃的形式是无政府主义。它认为，只要存在任何权力结构，即国家秩序，位置的分配从根本上就由主导/被主导这对关系所决定，而任何表象上的移动都使得本质的政治结构保持不变。

“革命”这个词本身变得可疑，因为它恰恰意味着矛盾的主要方面的改变。无产阶级党来占据国家位置，只是一种毫无意义的置换。他们更喜欢“颠覆”这个词，它指的是对任何主导位置的普遍瓦解，无论它是什么。他们将偏爱漂泊、漂流、无处可居，偏爱所有——至少在表面上——被排除在组合和置换游戏之外的东西。他们将宣称自己处于边缘：边缘人，比如《边缘报》（journal Marge）。

然而，这种对经典革命概念的起诉，远非最近几年政治上的发明，远非从68年五月风暴的腹部涌现出的新事物，它不过是施蒂纳的预言的一种既复杂又极其平淡的版本。

施蒂纳的学说将“反抗”与“革命”对立，其用词与源自68年五月风暴的小资产阶级革命运动的瓦解所散布的恶臭术语完全相同。唯一的区别在于一个微小的词汇变化，将“欲望”一词替换了圣马克斯（施蒂纳）所用的、更为坦率的“利己主义”一词。

至于其他方面，圣吉尔（德勒兹）、圣费利克斯（瓜塔里）、圣让-弗朗索瓦（利奥塔）在狂热的幻想大教堂中占据着同一个壁龛。他们认为“运动”是一种欲望的冲动，是一种流动的力量；认为任何机构都是偏执的，并且原则上与“运动”是异质的；认为没有什么能够对抗现有秩序，而是通过一种肯定的分裂而退出现有秩序；认为必须用自治——或联合，在某些地方对此有争论——来取代任何组织和任何丑陋的激进主义：所有这些大胆的修正主义，他们声称这些东西反对“极权主义”的马克思列宁主义，而高举边缘群众异议的闪耀新事物，这些都是马克思和恩格斯在《德意志意识形态》中——大约在1845年！——逐字逐句地拆解的东西，以清理出一条道路，以便对他们那个时代的革命实践进行最终连贯的系统化。

我们必须读了才能相信，这种思想的历史不变性是如此令人震惊，其主要论据正是它所谓的“根本性新意”。让我们来读一读：

\begin{quote}
“[反抗]是个人的起义，是一种上升，与由此产生的机构无关。革命旨在建立新的机构：反抗促使我们不再让自己被编入机构，而是相反地组织我们自己。

这不是一场反对现有秩序的斗争，因为如果它成功了，这个现有秩序就会自行崩溃，这只是我为了摆脱现有秩序而所做的努力。如果我退出现有秩序，它就死了，并进入解体。”\footnote{引自马克思、恩格斯，《德意志意识形态》，社会出版社，第414页。}
\end{quote}

现在，如果有人想要反等级制，更深层次地认为任何理性都会压制“运动”的纯粹自主性，认为革命组织建立在一种无法容忍的思想专制上，并且归根结底，有组织的激进主义及其机构等级制度，源于宗教心态？那么看看吧！圣马克斯就在那里：“直到今天，我们仍然服从于等级制度，并被那些依赖理念的人所压迫，而理念是神圣的。”\footnote{引自马克思、恩格斯，《德意志意识形态》，社会出版社，第198页。}

至于任何稍微一贯的革命计划都是一种“阉割”和对“力比多投入”的宗教性歪曲，让我们说，这是我们对列宁主义激进分子的“现代”和“激进”批判的陈词滥调，圣马克斯早在一百三十年前就在《唯一者及其所有物》的肮脏炉灶里烹饪过它，他将“天职-规定-使命-理想”这种神圣的普遍形式，与“自我愉悦”的规则完全对立。\footnote{马克思、恩格斯，《德意志意识形态》，社会出版社，第414页。}

马克思和恩格斯以无限的耐心拆解了这场概念闹剧的所有齿轮。他们在施蒂纳式“反抗”的狂喜颤抖的深处，以及其对任何有组织革命过程的“根本性新颖”对立（早在1840年代）中发现了什么？他们在将“颠覆现有事物”替换为不可言喻的“颠覆的存在”中无情地揭示了什么？\footnote{马克思、恩格斯，《德意志意识形态》，社会出版社，第414页。} ——就像我们“现代”的复读机们将“革命的欲望”与无产阶级政治计划的理性化和中心化领导对立一样——他们发现的是资产阶级、保守和法律主义的同一性原则的空虚。我就是我，因为我不是另一个；欲望就是欲望，因为它是我自己的。

\begin{quote}
“这种广受赞扬的‘独特性’，如此出色地与‘一致性’，与人的同一性区分开来，以至于桑乔（Sancho）将迄今为止存在过的个体都看作是一种物种的‘样本’，或者说差不了多少，但现在它却消失了，只剩下人与自身的同一性，正如警察所确立的那样，即一个人不是另一个人的简单事实。而桑乔，这个向世界发起进攻的巨人，现在被降级为护照办公室的职员。”\footnote{马克思、恩格斯，《德意志意识形态》，社会出版社，第485页。}
\end{quote}

无政府主义变成了警察，国家的颠覆变成了民事登记记录：我们又回到了结构主义，一种关于不变性、可置换原子、分类和组合的思想；它最终将运动简化为同一事物的换位。

无政府-欲望主义观念完全受到坚固的资产阶级客观主义的影响。它对事物的看法表面上与结构主义对立，因为它将所有结构的解体作为目标。但它从根本上是结构主义的，因为它只将矛盾视为一种替代运动，并排除了对两种统治类型之间任何质的差异的理解。

对于这种思想来说，任何分裂都是一种交换运动，任何斗争都是一种相互服务，任何胜利都是一种失败：在我让对手屈服于我的法则的那一刻，我就落入了它的法则之下，因为我占据了它曾经颁布法则的那个位置，即主导位置。

事实上，无政府主义不过是保守结构主义的简单镜像。漂流是组合的阴影。结构主义和欲望意识形态是根深蒂固的一对。它们远非对立，而是在它们与辩证法的共同矛盾中融合在一起。

事实上，辩证法，在其对主要方面理论的阐述中，同时思考了同一性（确实是在位置置换的范畴中）和非同一性（在作为整体的位置分配过程的质的突破范畴中），这遵循了其本身的分裂运动。辩证法使得结构与质、组合与力量的差异化评估之间的矛盾得以存在。

为此，辩证法必须分裂其“主要”的概念。

从第一层意义上讲，“主要”指的是一个过程的某个序列，正是从它出发，该序列得到了质的规定。
因此，在当前的历史序列中，法国确实是根据资产阶级的阶级主导地位而被定性为资本主义国家。这可以说是在过程中截取的一个静态切片。这是对一个结构性不变性的抽象规定，它固定了所考察序列的历史特殊性。主要矛盾的方面：资产阶级/无产阶级，在这里并非在其运动中的相互作用中被抓住，而是在其已确立的单一等级关系中被抓住，即根据主导/被主导这对关系。
因此，“主要”指的是对客观基础的唯物主义理解。它具有结构性价值。
更普遍地说，一个方面对另一个方面的统治的序列不变性，是唯物主义论点内在铭刻在辩证法论点中的方式。在我们的例子中，资产阶级的统治指向了资本主义的客观存在，指向了剩余价值的榨取过程，这构成了当今法国所有历史现象的一般物质基础，也是资产阶级国家权力的客观的、科学可分析的基础。
资产阶级是资产阶级/无产阶级矛盾的主要方面，在这里仅仅意味着，对社会进行唯物主义研究表明了其资本主义本质。
唯物主义铭刻在辩证法内容中的哲学浓缩表述为：在存在（物质）/思想（精神）的矛盾中，存在（物质）是主要方面。这可以说是否物主义论点的“辩证法”表述。

通过具体化这一普遍论断，我们发现了“主要”的一系列规定，它们构成了辩证法的唯物主义骨架。毛泽东有一段非常著名的论述，其全部意义就在于此：

\begin{quote}
“例如，在生产力与生产关系的矛盾中，生产力是主要方面；在理论与实践的矛盾中，实践是主要方面；在经济基础与上层建筑的矛盾中，经济基础是主要方面；各方面的位置不会相互转化。这种观点是机械唯物主义，而不是辩证唯物主义。

当然，生产力、实践和经济基础一般地说，是起着主要地、决定地作用的，谁否认这一点，谁就不是唯物主义者；但是必须承认，在一定条件下，生产关系、理论和上层建筑，又可以转过来起主要地、决定地作用。”\footnote{毛泽东，《矛盾论》，《毛泽东著作选集》第一卷，人民出版社，1991年，第323-324页。}
\end{quote}

总之：
a) 为了成为唯物主义者，必须承认一系列方面（实践、生产力、经济基础）“一般地说”占据主导地位，它们是其与对立方面（分别是：理论、生产关系、上层建筑）统一的矛盾的主要方面。
因此，主要方面在某种程度上具有固定性，这确实是唯物主义锚定在某些辩证法内容中的方式。

b) 为了成为辩证法家（即不成为机械主义者），也必须承认这种固定性的否定。
如果说主要方面在战略上（“一般地说”）的固定性证明了唯物主义在辩证法中的存在，那么其在战术上（“在一定条件下”）的非固定性则证明了辩证法在唯物主义中的存在。

从所考虑的例子来看，可以说：唯物主义是通过战略性地固定矛盾各方面的位置来构建矛盾的东西；辩证法是通过思考位置的颠倒，即各方面分配的非固定性来反驳结构的东西。

“主要”一词的第一种含义，指的是稳定的主导现象，实际上主要是其唯物主义的含义。它必须被辩证化。

事实上，如果说辩证法首先是矛盾运动的规律，那么我们不能满足于简单的结构性规定，因为它本身对这种运动只字未提。法国是一个资本主义国家，资产阶级是统治阶级，这是对的。但历史的动力不是资本主义体系，而是阶级斗争。

“阶级斗争”中当然包含了经济基础、下层建筑，它们决定了阶级本身的存在，以及它们之间暂时的从属体系。但“阶级斗争”首先指的是对抗的过程，通过这个过程，推翻现有统治的条件得以累积，也就是被主导的方面成为主要方面的过程。对“主要”的完全辩证理解因此不仅必须理解事物的状态，还要理解其趋势。

从事物状态和结构的角度来看是主要的东西，从趋势的角度来看可能变成次要的。因此，毛泽东在1970年5月的声明中毫不犹豫地说，在一个仍然被帝国主义和社会帝国主义统治的世界里，“主要的趋势是革命。”\footnote{毛泽东，《全世界人民团结起来，打败美国侵略者及其一切走狗！》，北京出版社，1970年，第2页。}

帝国主义和社会帝国主义从结构上是主导的：世界上大多数人民仍然遭受其压迫。

但从历史运动的角度来看，它们已经不再是主要的了。正如毛泽东所说，辩证法在于从未来的角度看待事物。要在当下本身从未来的角度看待事物，就必须抓住这个当下作为一种趋势，作为力量的增长或衰退，作为积累，作为突破，而不仅仅是作为状态或形态。

最终，“主要”指的是当一个矛盾的方面在其与另一个方面的关系中被理解时，它所具有的双重性质。“主要”遵循了分裂过程的分裂思想。每个方面都同时处于一种给定的主导或依赖关系中，并处于对这种状态的否定中。

世界人民在大多数情况下都遭受着资本主义和帝国主义的压迫，但主要的历史趋势是“国家要独立，民族要解放，人民要革命。”而这种意愿是反抗反动派，这种意愿是正确性。人民是被主导的，帝国主义是主要方面。但人民，正因为他们是被主导的，他们真正的存在，他们的历史存在，就是摧毁帝国主义；而这种历史存在，在世界舞台上，是主要力量。“主要”分裂了，“主要”变成了它的对立面。

每一个方面，正是因为它只有在参与分裂过程中才有存在，所以它本身被拉扯在它对作为过程（即作为统一）的分裂的质的从属，以及它对这种质本身的转化运动之间，而这种转化运动的动力是两个方面之间不间断的斗争，以及它们之间关系的持续改变。

结构的存在是一种等级组合，但它的存在，即它的历史，与其自身的毁灭相融合。结构除了其自身灭亡的运动之外，没有其他存在，而矛盾的每个方面都通过其分裂为“为结构而存在”和“为结构解体而存在”来反映这种过渡性的存在方式。无产阶级将在生产过程中被结构性地定义为被剥削阶级，并被反结构性地定义为承载着消灭资本主义生产方式的革命阶级。

在第一种规定中，无产阶级本质上是被主导的，在第二种规定中，主要方面不是它的被主导存在，而是它自身统治得到肯定的历史趋势。从作为整体的资本主义社会的角度来看，说资产阶级是矛盾的主要方面和无产阶级也是主要方面，都是正确的，这取决于我们在“主要”的分裂中，将结构还是趋势放在主导位置。

辩证法，可以说，它本身就是辩证的，因为其反映现实的概念性操作者，也都同样是分裂的。这就是列宁所说的“概念的多方面的、普遍的灵活性，这种灵活性甚至达到了对立面的同一性”\footnote{列宁，《哲学笔记》，社会出版社，第108页。}。

一个辩证法概念的真理，永远不在于其某一种单一的含义，而在于其两种对立含义的灵活的矛盾运动。因此，“主要”的完整概念将是结构和趋势的矛盾统一。单方面考虑其中之一，就会陷入形而上学。

这种形而上学反过来又是“极左”或修正主义政治立场的浓缩。

如果确实忽略了结构性元素，就会把趋势当作已经实现的事实。因此，盖斯马尔在68年五月风暴中解读出武装内战的迫近性，直接导致了前“无产阶级左翼”的毁灭性唯心主义冒险。因此，《红色人道报》篡改了中国关于欧洲战争不可避免的论点（这是一种趋势性的断言），表现得好像苏联侵略已经发生，并与最糟糕的反动派——例如“新法兰西行动”——进行了可笑和可耻的“爱国统一战线”操作。在这两种情况下，一种趋势都被当作了既成事实。趋势上的主要方面与结构上的主要方面融合了。

如果忽略了趋势性元素，就会不可避免地以旧事物的名义压制新事物，支持现有秩序。这会陷入机会主义的等待。没有必要去辨别和巩固趋势，结构会满足我们的需求。这是第二国际修正主义者的学说：反动的经济主义，它假设，资本主义下层建筑通过其自身的客观运动就会导向社会主义。只有结构决定未来，只有它才是趋势。结构上的主要方面与趋势上的主要方面融合了。

这些偏离的形而上学都回到了“合二而一”的原则。它们放弃了“主要”概念的辩证分裂，并对称地为“极左”冒险和右翼机会主义背书。

真正的辩证研究总是同时操作对主要方面的结构性规定和趋势性规定，这是一种分裂的关联。

事实上，一切都归结于此：在给定的状态下，或在其发展的特定序列中，任何矛盾都会为其方面分配一个确定的位置，这个位置本身由其与另一方的位置关系来定义。主要方面就是在主导和被主导的位置分配中被抓住的。从这个意义上说，辩证法是一种位置逻辑。

相反，当在它不间断的阶段性运动中被抓住时，任何矛盾都使力量相互对抗，而这些力量的性质是差异性的：在评估力量时，从矛盾运动的角度来看重要的，不再是其暂时的从属或主导状态，而是其增长或衰退。在其趋势性或纯粹历史的一面，辩证法是一种力量逻辑。

因此，核心的辩证问题是：位置逻辑和力量逻辑如何衔接，而不融合？

\section{力量与位置}

\subsection{矛盾的解决。什么会消亡？}

红色高棉占领金边：一个历史序列结束了，因为一个矛盾被解决了。

矛盾的解决，不是，也永远不是其各方面的综合。我们在这里掌握了马克思主义最稳定和最古老的真理之一：斗争摧毁和改造，它不是一场在“起源”的姿态中策划的和解的影子戏。从《宣言》开始，马克思和恩格斯就毫不含糊地肯定了这一点：

\begin{quote}
“[…]压迫者和被压迫者，始终处于相互对立的状态，进行着不间断的、有时公开、有时隐蔽的斗争，每一次斗争都以整个社会革命改造或者斗争的各阶级同归于尽而告终。”\footnote{马克思，恩格斯，《共产党宣言》，《马克思恩格斯全集》第一卷，人民出版社，1995年，第276页。}
\end{quote}

如果矛盾的一方不可逆转地掌握了主导位置，那么就会出现整体的重铸，即过程的质的突破；如果力量的对称性不允许位置的改变，那么过程本身就会在两方的相互消灭中结束。

一个矛盾的解决要求某物必须消失。唯物主义辩证法不像黑格尔辩证法那样，是一种关于“在死亡中保持自身”的思想。\footnote{黑格尔，《精神现象学》，奥比尔出版社，第一卷，第29页。}或者更确切地说，辩证法将死亡一分为二：确实有在新的冲突性条件下仍然存在的东西：例如，资产阶级在无产阶级专政下仍然存在，以权利持续存在的客观形式，以其渗透到党内的主观形式。

但也有完全消亡的东西，不留任何痕迹的东西。只有能够将对新事物的认识推向其不可避免的另一面：旧事物必须消亡，才是真正的革命思想。

唯物主义辩证法像黑格尔一样肯定，真正的生命不是那种“在死亡面前退缩，并保持自己纯粹不被摧毁”的生命。但它对死亡的分裂思想更为强大。黑格尔想要“牢牢抓住死者”，正是因为一切都融入了绝对的最终循环，所以精神的生命无需惧怕否定。

没有什么会永远失去。相比之下，唯物主义辩证法直面损失，直面一去不复返的消失。激进的新事物之所以存在，是因为有永远不会被任何审判号角唤醒的尸体。

在文化大革命最激烈的时候，中国有句说法：修正主义的本质是对死亡的恐惧。这句话分裂为两部分，其唯心主义和主观的一面——林彪的一面——以及其真理的一面：修正主义者无法承受资产阶级的死亡。

修正主义是资产阶级生存的最终化身，是那个必须一劳永逸地离开位置的人的狡猾和伪装，他为了生存而穿上其对抗者的旧衣服：伪装成社会主义的资产阶级专政，贴上假马克思主义标签的资产阶级思想。死者只有在被误认为是杀死他的人时，才能获得苟延残喘。

我们说，矛盾的解决包括死亡的一部分。为了新整体的出现，为了不同过程的出现，为了另一个统一的分裂，现实的一个碎片作为运动的废物而坠落。这就是“历史的垃圾堆”这个战斗性表述的全部合理性。解决就是抛弃。历史工作得越好，它的垃圾堆就越满。

黑格尔用“Aufhebung”（扬弃）来指代从旧到新的过渡：取消并保留。我们可以翻译为“接替”。另一个事物，却又存在于连续性中。消亡的事物在其对立面中实现自身。“扬弃”对我们来说，必须被分裂：无产阶级确实是资产阶级的接替者，但不是以工厂下班后夜班工人接替白班工人的方式，甚至也不是以世代接替的方式。

因为资产阶级作为剥削者历史的最终继承者，最终将一无所有。而毛主义者将接替法国共产党在工人阶级中的位置，不是通过托洛茨基主义者那种批判性和寄生性的陪伴，而是通过对修正主义者领导存在的物质性摧毁。

我们必须用辩证法来思考，不仅是死亡，还有骨灰的散落。

正是通过分裂死亡，马克思列宁主义者才能对其进行评估。在那些持续存在和那些被废除的东西中，什么是主要的？巴黎公社的失败，是对法国无产阶级一整代领导人的不可挽回的损失。这种损失是如此强烈和真实，以至于我们至今仍在承受。接替，是对公社的总结，由马克思进行系统化，由列宁和俄国无产阶级以工人胜利的方式实践。公社社员的死亡，其主要方面是胜利的无产阶级革命的活生生的未来。马克思在当时就看到了这一点：

\begin{quote}
“工人阶级的巴黎及其公社，将永远作为新社会的荣耀先驱而受到赞扬。”\footnote{马克思，《法兰西内战》，作品全集，第二卷，第256页。}
\end{quote}

反过来，例如，殖民体系的死亡，随着葡萄牙“帝国”的崩溃（几内亚-佛得角、莫桑比克、安哥拉）几乎完成，并没有预示着任何内在于该体系的事物。

没有任何殖民“价值”会在未来找到其概念的荣耀。这个被摧毁的、不可挽回的、不孕的体系：反抗人民的力量将其分解并注定其进入垃圾堆。

一个如此彻底被消灭的方面，在其发展的整体性中被触及：当它的起源本身也在对抗中被取消时，它就结束了。这就是为什么毛泽东认为美国黑人的反抗是整个帝国主义体系彻底崩溃的一个决定性时刻：

\begin{quote}
“从压迫和贩卖黑人中兴盛起来的万恶的殖民主义和帝国主义制度，必将随着黑人的彻底解放而告终。”\footnote{毛泽东，《支持美国黑人反对美帝国主义种族歧视的正义斗争的声明》，北京，1963年。}
\end{quote}

当一个方面在其初期所反对的增长，反过来又获得胜利时，这个方面就离其消亡不远了。

因此，在一个正在走向终结的现实中，分离并有序地存在着：融入未来的东西，在新对立面统一的规律下进行转变的东西，以及作为这种新事物的损失和瓦解的反面而存在的东西。所有消亡的东西都在矛盾中被接替和被剔除。根据接替还是废物占据主导，死亡在辩证法上是肯定的还是无效的。

\begin{quote}
“中国古代的作家司马迁说过：‘人固有一死，或重于泰山，或轻于鸿毛。’为人民利益而死，就比泰山还重；替法西斯卖力，替剥削人民和压迫人民的人去死，就比鸿毛还轻。”\footnote{毛泽东，《为人民服务》，《毛泽东著作选集》第三卷，人民出版社，1991年，第1004页。}
\end{quote}

柬埔寨傀儡政权的覆灭，属于“鸿毛”和废物的范畴。朗诺集团从历史舞台上消失了，革命内战以人民力量的彻底胜利而告终。会有人说红色高棉来占据了傀儡曾经的位置吗？从某种意义上说，是的，因为一个国家被另一个国家所取代，一个专政被另一个专政所取代。

从形式上看，一个矛盾过程的结束和解决在于不对称性的翻转，这是正确的。权力易手，次要的变成了主要的，过程的整体质的规定不再由同一个主导方面来支配。

对过程本身的称谓也必须改变：革命内战阶段之后是人民民主专政阶段。国家，即所有统治现象的集中形式，也不再具有相同的名称。由买办和受帝国主义豢养的军人组成的傀儡国家，被一个建立在工人和农民联盟基础上的国家所取代。正如《国际歌》所唱的：“世界换新颜”，社会力量交换了位置，对“谁压迫谁？”这个问题的回答被颠倒了。

\subsection{矛盾的解决。什么会改变？}

然而，所有这一切都描述了辩证序列中一个突破的形式，但没有触及其内容。人民民主专政国家无疑是傀儡国家的颠覆。在一个阶级关系组合的理解中，我们甚至可以说它实现了这种颠倒：广大的工农群众对旧独裁者实行了专政。

然而，这种颠倒的历史发展是一种激进的新事物，单靠颠倒的思想是无法理解的。这里甚至存在一种无政府主义反对意见的风险，它可以声称是来自马克思的：即使是颠倒的统治关系的存在，难道不本身就构成了一种本质的意识形态不变性的核心吗？在《宣言》中，马克思确实从剥削关系的一般持续性中，推断出意识形式的稳定性：

\begin{quote}
“[…]过去的一切社会，尽管形式各不相同，都有一个共同的特征，就是社会的一部分人对另一部分人的剥削。所以，毫不奇怪，过去各个世纪的社会意识，尽管经过各种各样的变动，总是在某些共同的形式中运动的，这些意识形式，只有当阶级对抗完全消失的时候，才会完全消失。”\footnote{马克思，恩格斯，《共产党宣言》，《马克思恩格斯全集》第一卷，人民出版社，1995年，第276页。}
\end{quote}

在这里，让我们仔细跟随形式/内容辩证法的曲折。一个不变的矛盾内容（阶级对抗）的存在，使得其形式（不同的剥削方式）的可变性，在主观因素（社会意识）中留下了一个形式上的不变性（共同形式）。（这种形式上的不变性反过来又分裂为所有剥削者共同的意识形式，和所有反抗的被剥削者共同的意识形式：这是所有伟大人民反抗中共产主义意识形态不变性的问题。我们将在《意识形态、反抗、理论》这本小册子中讨论它。）

如果我们坚持将矛盾的解决仅仅视为位置的颠倒，我们可能会说：一个不变的矛盾内容（一个主导/被主导位置体系的存在），无论其形式如何变化（国家革命），都会导致意识形式的固定性（没有思想革命）。在这一点上，客观辩证法会麻痹主观辩证法。更糟糕的是：剥削体系的颠覆将与对压迫实践的决裂完全脱离。

这是因为位置的互换只包含对辩证分裂运动的思考，而且仅仅是从两个方面的关系来看。每个方面的存在都在于它与另一个方面（另一个位置）的关系。每个方面的力量都在力量关系中耗尽。占据对手的国家位置确实证明了两个阵营之间力量关系冲突性转化达到了终点。但这种占据没有说明参与这种关系的力量的内在性质。

甚至说人民阵营现在拥有比反动阵营更优越的力量也是不够的。这种优越性的肯定只不过是位置改变这一显而易见的事实的重申。如果红色高棉能够扫除他们的对手，那他们无疑是最强的。在这个调查阶段，力量被理解为一种外部的、简单的量化关系。

但是，在力量的量化概念中，除了位置逻辑所提供的东西之外，没有更多东西。衔接力量逻辑和位置逻辑，不能仅仅归结为确定位置改变的力量条件。事实上，如果力量的比较仅仅局限于量化的思维，那么位置的改变不一定意味着过程的质的转化。

我们说过：质的突变需要回答这个问题：什么会消亡？现实中堕落的部分在哪里？力量的量化思想无法提供这个答案。它表面上告诉我们什么在变化（位置），却对什么会消亡（因此什么会诞生）漠不关心。修正主义的选举承诺遵循了这条道路：“权力”将被左翼联盟所取代，共同纲领将扭转局面。

但是什么必须消亡呢？听他们说，什么也不需要。也没有什么会诞生，除了已经存在的：党、工会、政府。这个没有损失的历史，这条没有垃圾堆的大道，引发了辩证法的怀疑。他们在对我们撒谎，要么是不想改变任何东西，厌恶任何历史性的诞生；要么是隐瞒了代价、损失和死亡。他们确实在隐瞒：因为损失是属于群众的，群众在承诺的实践真理中，被迫承受修正主义的无法忍受的独裁，承受致命的社会法西斯主义。

\subsection{次要过程：当什么都没有改变时，消亡的是人}

辩证法抓住分裂、致命的、正在诞生和肯定的事物。凡是位置互换单独主宰现象的地方，改变就变成了它的对立面：持续性，以及它的反面：屠杀。

让我们来看一个标题对称的现象：一场帝国主义间的战争，帝国主义对抗帝国主义，而不是红色高棉对抗傀儡。关于位置，关于力量，该说什么？在1914年，法国帝国主义和德国帝国主义的战争目标是占据对手所占据的位置。德国想扩大其在欧洲的影响范围，并全面参与对非洲的帝国主义掠夺。

法国想与英国一起保卫其半殖民霸权，并夺回阿尔萨斯-洛林。因此，两个帝国主义中任何一个的胜利，都将以位置的颠倒而告终。导致这些颠倒的无疑是两个帝国主义之间的力量关系。但正是这些力量本质上是同一种性质：帝国主义力量。

因此，矛盾在位置的片面规律下运动。对力量的评估在量化的思维中展开，在同质的元素中展开。归根结底，拥有最多的人、最多的武器和最多的盟友的一方，不可避免地会获胜。这并不是说量化和武器是战争的决定性因素：我们知道，自克劳塞维茨以来，“道德力量”和群众的支持是不可战胜的唯一因素。然而，这需要存在不对称性：克劳塞维茨以辩证法家的身份思考战争，有一方是进攻者，有一方是防守者。

防御的战略优势在很大程度上来自于人民群众的充分动员。然而，帝国主义间战争总体上是反革命的，其目标外在于任何人民的愿望，除了掠夺和侵略之外，没有什么可捍卫的。

因此，不奇怪，在帝国主义间战争的核心，同一性压倒了差异性，因为历史运动的决定性参与者（人民群众）作为矛盾的一极，被排除在整个战争过程之外。

帝国主义间战争，无论其多么野蛮和绝对，都不过是现实矛盾运动的一种退化、削弱、最低限度的形式。马克思一百年前就以完美的清晰度说过：

\begin{quote}
“旧社会所能达到的最高英雄主义，就是一场民族战争；而现在已经证明，这不过是政府为了推迟阶级斗争而进行的纯粹的欺骗，一旦阶级斗争爆发为内战，这种欺骗就会被抛弃。阶级统治再也无法隐藏在民族制服之下，民族政府在对抗无产阶级时是合为一体的！”\footnote{马克思，《法兰西内战》，作品全集，第二卷，第256页。}
\end{quote}

因此，帝国主义间战争与其说是一场战争，不如说是一种欺骗，因为其内在规律使得在对抗的核心，统一压倒了分裂。这并非一个辩证法的例外。如同任何概念一样，对抗的概念也分裂了：某些对抗性矛盾中存在非对抗性元素，而这个元素的运动决定了在任何公开对抗的过程中，无论其多么暴力，统一和斗争之间的关系：

\begin{quote}
“随着事物的具体发展，某些原来是非对抗性的矛盾，发展成为对抗性的矛盾，而另外一些原来是对抗性的矛盾，则发展成为非对抗性的矛盾。”\footnote{毛泽东，《矛盾论》，《毛泽东著作选集》第一卷，人民出版社，1991年，第323-324页。}
\end{quote}

中国同志正是在这个对抗性分裂的过程中，才能够在无产阶级文化大革命期间，提出了用解决非对抗性矛盾的方法来处理某些对抗性矛盾的指示，否则这是无法理解的：

\begin{quote}
“以解决人民内部矛盾的方式来解决这些敌我之间的矛盾，有利于加强无产阶级专政，瓦解敌人。”\footnote{《在中国共产党第九次全国代表大会上的报告》。}
\end{quote}

帝国主义间战争表面上将对抗推向了极致；然而，其内在的非对抗性主宰了其存在形式，即累积而非创造，无止境，血腥，且极其不孕。（对抗性自身内部的分裂及其对立面，规定了分析我们当前局势的一个关键点：两个超级大国——苏联和美国——的关系，受到竞争和勾结的严密辩证法的支配。）

这就是为什么它脱离了克劳塞维茨的范畴。这也是其风格的原因：一场愚蠢的相互屠杀，对物资的无谓挥霍，一种漫长而痛苦的力量消耗，它们之间的联系越紧密，就越是本质上没有区别。

这种杀戮的对称性，这种类型的过程是矛盾的弱结构形式。量化的思想在这里找到了它的养料。对力量的同质评估与位置的互换（胜利或失败）之间存在一种机械关系。因果关系是线性的。位置逻辑从属于力量逻辑。让我们给精神分析一个花（修辞的花）：我们将这种类型的矛盾称为次要过程。

它确实是次要的，因为它是在另一个矛盾，一个主要的矛盾中产生的游戏，在这个例子中，是帝国主义与世界人民之间的矛盾，它使得帝国主义变得对称，并使得它们之间的野蛮搏斗成为一种现象，这种现象本身除了重新分配无关紧要的位置之外，不带来任何东西。

我们会说次要过程的解决没有损失吗？当然不是。损失以一种完全消极的形式出现：人口的屠杀。在14-18年的大屠杀中死去的数百万工人和农民，是帝国主义毫无意义的位移的废物。那么，要去哪里寻找必须消亡的东西呢？除了那些被排除在战争动机和利害关系之外的人们之外，还能去哪里呢？

我们说过，次要过程，作为一个整体，在全球范围内与它从其位置体系中排除的东西进入了主要矛盾：人民客观地（迟早也会主观地；看法国国内，1917年的反抗）反对帝国主义战争。正是在其主要对立面中，次要过程找到了必须被剔除的方面，以便其解决能够到来。

（这是一条普遍的辩证法规律，我们在这里只考虑其退化的运作方式。它也同样适用于积极的方面。为了完全解决——例如——工人与农民之间的矛盾，必须消除旧社会遗留下来的三大差别（脑力劳动与体力劳动、城市与乡村、农业与工业）。

然而，反对这种消除的历史力量，正是党内和国家中仍然存在的修正主义资产阶级，其主要客观支持恰恰在于在漫长的过渡时期中这些差别的必然持续。因此，我们可以完美地说：工人与农民之间的（主要是非对抗性的）矛盾定义了一个次要过程，而其主要过程是围绕政治权力和国家问题的无产阶级与资产阶级之间的矛盾。因此，为了解决次要过程，必须消灭渗透到工人阶级、党和国家中的修正主义资产阶级，这需要革命的方法：一系列的文化大革命。修正主义者因此是通向共产主义的道路上的废物，是死去的方面。我们将在专门讨论不同类型矛盾的小册子中重新讨论所有这些问题。）

在14-18年死去的法国或德国工人和农民，正是为了他们的民族帝国主义而死的，恰恰是因为他们构成了整个帝国主义体系的对抗性方面。

所有本质上对称的过程都通过位置的改变来解决，其代价是从异质力量中获取。对称性是以不对称性为代价的。位置是通过阉割力量的差异性来获得的。

然而，如果次要过程是倒退的，那么它又被重新启动了：法国或德国的农民和工人也为俄国无产阶级革命的到来而牺牲。作为死去的方面，辩证地唤醒了人民真正的生命，即加速了主要过程。

对战争的反抗，使得无产阶级和人民群众以其本质的同一性而得到确定：无产阶级革命对抗帝国主义野蛮主义。因此，对帝国主义间战争的对称性和无关紧要性的认识，打破了资产阶级民族主义的堤坝——这种虚假质的差异性（我们法国人，真的和德国鬼子不一样），旨在让人民占据帝国主义者所要求的那个位置：死者的位置。

在资产阶级民族主义的破裂中，阶级意识将帝国主义间战争视为一种欺骗和统一的置换。它整体上以布尔什维克的革命失败主义口号来对抗它。主要过程战胜了次要过程。

但在主要过程中，力量和位置的衔接方式完全不同。

\subsection{主要过程：力量的闪耀}

主要过程本质上是不对称的：人民的力量和帝国主义的手段无法通过量化的思维来评估。在红色高棉与傀儡之间的胜利矛盾中，在场力量的异质性使得它们严格来说没有任何共同的衡量标准。资产阶级思想对大帝国主义在小反抗民族面前的失败，总是感到永久的震惊。

这是因为，如果坚持将力量关系的概念严格从属于位置逻辑，就会打开一个纯粹量化的思想空间，在这个空间里，显然美国帝国主义——例如——是最强的。在人员、武器、火力、摧毁能力上，其中一个对抗者拥有压倒性的优势。然而，正是它被打败了。正是它必须让位。我们可以清楚地看到，在这种情况下，对力量的评估问题不能通过量化计算的对抗来解决，因为这只会机械地导致位置的互换。

让我们思考一下毛泽东的这个政治论断：

\begin{quote}
“一个弱国能够战胜一个强国，一个小国能够战胜一个大国。小国人民如果敢于站起来斗争、拿起武器、掌握自己国家的命运，就一定能够战胜大国的侵略。”\footnote{毛泽东，《全世界人民团结起来，打败美国侵略者及其一切走狗！》，北京出版社，1970年，第2页。}
\end{quote}

在同质的逻辑中，这个论断是无法理解的。它实际上意味着，在某些只涉及其内在规定的条件下，最弱的却比最强的更强。

然而，这是一个本质性的辩证论断，一个真正的马克思列宁主义指令。列宁早已强调，起义本身，作为一种胜利的力量行动，并没有结束无产阶级的劣势：它证明自己是最强的，但它仍然是最弱的，并且仍然必须赢得真正的优势。在这种情况下，对国家位置的夺取，远非机械地解决了力量关系问题，而是重新引发并加剧了它：

\begin{quote}
“无产阶级专政是新生阶级对一个更强大的敌人，对被推翻后反抗力增加了十倍的资产阶级进行的最英勇、最无情的战争。”\footnote{列宁，《共产主义运动中的“左派”幼稚病》，作品全集，第三卷。}
\end{quote}

从位置上来看最弱的（小国，被统治的无产阶级）可以根据力量（但那是一种不同于由位置规定的量化力量）而战胜，并颠倒位置。然而，它在很长一段时间内仍然是最弱的，而且正因为如此，它现在还必须保卫所夺取的位置，从而锻造它以前没有的东西：位置上的优势。

所有这些辩证法都分裂了力量的概念，就像分裂了弱点的概念一样，并交织了两种评估力量关系的标准。对于我们毛主义者来说，作为革命未来的坚定承载者，这是一种决定性的思想，因为它指出了我们——为数不多的人——如何注定要战胜双头资产阶级怪物：经典资产阶级和新修正主义资产阶级，尽管我们在很长一段时间内仍然是最弱的。

为了理解最弱者的力量优势，为了理解红色高棉同时战胜美国帝国主义和社会帝国主义愿望的胜利，必须引入一种矛盾力量之间不可比拟的观念；一种异质评估体系的观念，这使得评估力量关系的标准本身是分裂的。这是抓住一个主要过程的一个本质特征。

例如，马克思列宁主义在著名的“人民战争是不可战胜的”这一公式中综合了这一点。这个公式实际上意味着，人民战争所使用的力量类型无法与对抗力量所拥有的武器库相比较。说一种力量是不可战胜的，就是赋予它一种内在的、无法与其他力量相比较的性质。

人民战争的伟大原则：政治第一性而不是军事，群众的支持是主要因素，人民民主革命纲领，无产阶级的领导等等，解释了在过程的一极，一个方面是如何形成的，它的异质性本身保证了其必然的胜利增长。因此，矛盾思想不再是简单地通过力量的结构性（组合性）评估来复制位置体系。过程本质上是不对称的，是一种不可简化的不对称性，因为它具有质的性质。

这种不对称性不能再根据那种构成对位置进行量化评估的逻辑扩展来思考。有必要将力量的本质作为一种肯定性的品质引入。辩证法在这里需要的，是完全不同于关联原则的东西。它要求抓住分裂过程中的一种内在的质的非连续性。

反过来，力量的质的异质性颠覆了位置互换过程的内在性质。一旦最初被主导的方面的胜利增长与其不可还原的质的差异性得到肯定，这个方面就不能严格地占据旧主导方面的位置。质的新颖性必然延伸到统治本身的内容。新方面占据主导位置的方式，与其旧的占据方式是不可还原的。

一场帝国主义间战争的结果是位置的组合性重新分配；一场革命战争的结果是对位置体系本身的重铸。傀儡的统治与红色高棉逐步建立的人民民主专政国家之间没有任何共同之处。因此，我们必须说，这次是位置逻辑从属于力量逻辑。历史的新事物，在人民阵营的加强中得到了肯定性的增长，它带来了一种不同的规则，以及不同类型的位置关联。

一个次要过程使其定义的分裂对称化，置换位置，并给异质性分配死者的位置。一个主要过程则最大程度地使其分裂不对称化，重铸位置体系，并给旧的位置占据方式分配死者的位置。

最主要的过程，就是无产阶级革命。马克思主义的一个本质论点是，无产阶级不可能利用资产阶级国家机器。无产阶级革命只有在有效摧毁这个国家并建立根本上新的权力形式时才可能被构想。这等于是说，无产阶级不能简单地从主导/被主导这对关系的互换角度来肯定其革命同一性。

严格来说，无产阶级不能占据资产阶级的位置，它必须摧毁这个位置本身。列宁已经预见并抨击了纯粹结构性的、可置换的革命观念所造成的破坏：

\begin{quote}
“那些从小业主的角度看待战胜资本家的人们：‘他们发财了，现在轮到我们了！’，正在产生新一代的资产阶级。”\footnote{列宁，《在全俄中央执行委员会会议上的讲话，1918年4月29日》，《列宁全集》，第二十七卷，第312页。}
\end{quote}

我们不要将这个公式限制在其经济含义上。在一切方面，只要不将其次要化，即不背叛它，就永远不涉及以“现在轮到我们了”的方式实践主要过程。因为这个“轮到我们”预设了那个我现在轮流占据的位置的不变性。

在次要过程中就是如此：在1914年，对于喀麦隆人或摩洛哥人来说，要忍受法国人还是德国人来占据殖民地位置，这无关紧要。但在主要过程中，位置本身必须被改变，在一种无法比拟的、异质的、无法被旧主导力量所占据的地方的力量的影响下。

因此，无产阶级行使其统治的地方（无产阶级专政的国家），其本质不是持续存在，而是消亡。无产阶级占据主导位置的过程，同时也是这个位置解体的过程。

在这里，力量和位置概念的双重分裂清晰地衔接起来：国家无产阶级性质的标准是其自身消亡的真实矛盾运动。主导位置向群众本身移动，没有独立的机构，这证明了无产阶级作为一种不同质的力量，其内在本质得以持续和发展。

同时，在很长一段时间内，国家仍然是分离的。因此，一场激烈的阶级斗争爆发了，其最终赌注是，尽管客观上存在将过程次要化、固定位置、使力量对称化的力量，但如何保持过程的主要性，这些力量以冲突的方式从四面八方渗透进统治阶级（无产阶级）、党和国家。

在特定时刻对无产阶级阶级力量的评估，必须从其内部产生的质的新事物的角度来看，从在实践中构成异质力量的东西来看。任何表面上是工人的力量，但显然与资产阶级力量有着共同衡量标准，都没有机会融入一个主要过程。

在这一点上，与修正主义者，或受其影响的人，之间的矛盾，特别是在评估力量的方式上爆发了。对于修正主义者及其追随者来说，力量实际上被简化为位置。修正主义是一种资产阶级小集团的思想，它与另一个小集团（马歇、工会和国有化，对抗吉斯卡尔、政府机构和经典垄断）进行竞争，它只想理解次要过程。任何出现的主要新事物，它的本能反应就是将其压垮：德国犹太人！极左分子！冒险家！等等。

修正主义政治由对称性主宰：它从选举（次要性的典型思想，在本质不变性中有规律的交替）展开，通过政变，或入侵者的苏联货车，到反革命内战。

在所有情况下，都是资产阶级力量对抗资产阶级力量，以占据不变的国家。对于修正主义者来说，强大，就是简单地能够占据那些目前由资产阶级其他派系所占据的位置。

因此，在葡萄牙，修正主义者在4月25日之后，着手系统地投资于国家机器的整个文职部分（市政厅、行政机构、工会、媒体、教育等）。这种投资是在脱离群众真实历史运动的情况下，而且通常是与其对抗而进行的。它效仿了从卡埃塔诺反动政权继承下来的权力体系。

在最严格的意义上，对位置的占据是修正主义在权力问题上的政治学说。即使在修正主义者完全缺席、与群众没有任何联系的地方，他们也认为通过官僚主义地殖民那些空缺的权力职位，可以对人民群众建立他们的独裁。因此，在北部，他们在蔑视的农民面前，将自己塑造成一种新的领主，拥有权力和特权，这些特权比经典的反动权力更难以忍受，因为它们没有任何活生生的传统根基。

我们在这里看到了修正主义对历史过程思想的赤裸裸的扭曲：对他们来说，一切都是次要的。一切都组织着对异质性运动的无知和压制，一切都将历史简化为位置的组合，以及唯一的量化力量评估平面。葡萄牙农民群众对修正主义的强大反抗，尽管考虑到法西斯主义代理人的狂热活动而显得模棱两可，但其本质动力正是那种修正主义思想和实践试图排除在历史舞台之外的东西的涌现：人民力量无限的、肯定的能量，它无法被任何规范、任何组合所还原。

葡萄牙北部的农民群众有力地拒绝了在潜在的社会法西斯政变的次要过程中，扮演被轻视的死者角色。这再次证明了任何承载未来的力量最终都是不可战胜的，因为它以与现有位置分配体系在质上不连续的内在本质，矛盾地展开了自身。

关于阶级矛盾的本质是不对称和质的，以及关于主要过程和次要过程之间关系的，现代修正主义者和马克思列宁主义毛主义者之间的绝对对立，在六十年代围绕战争与和平的问题上以典范的方式爆发了。

当时在已故赫鲁晓夫领导下的苏共说什么了呢？它非常精确地说，原子弹使得所有与战争有关的矛盾都对称化了；从今以后，任何涉及核升级风险的战争，而核武器以（恐怖）平衡作为其发展形式，战争已经成为一种本质上是次要的过程（即不惜一切代价要避免的）。正如苏共中央委员会天真地宣称：“原子弹不遵守任何阶级原则。”\footnote{《苏共中央致党组织和全苏共产党员的公开信》，1963年7月。}

修正主义者完全认同资产阶级的“理论”，即核武器的出现从根本上颠覆了战争的内在原则。对他们来说，技术是历史的动力，而不是阶级斗争：“在本世纪中叶发展起来的核导弹，改变了人们对战争的看法。”\footnote{《苏共中央致党组织和全苏共产党员的公开信》，1963年7月。}

这种“看法的改变”集中在对辩证法的摒弃上：对于修正主义者来说，战争的概念是不可分割的。中国同志立即对此提出了反对：如果战争本质上是一个对称的过程，如果“死者的一部分”处处都永远由人民占据着（“原子弹不分谁是帝国主义者，谁是劳动者，它打击的是整个区域；为了一个垄断者，就会消灭数百万工人”\footnote{《苏共中央致党组织和全苏共产党员的公开信》，1963年7月。}），那么马克思主义者与小资产阶级和平主义者之间一直存在的战略区别——区分正义战争（属于一个主要过程，其中旧的剥削体系占据了死者的位置）和非正义战争（确实是对称的，对人民是致命的）——又何在？

战争概念的不可分割性——它归因于技术的“客观性”——强制性地将力量与位置进行了根本性的等同：被核武器的共同量化标准所制约的“社会主义”阵营和帝国主义阵营，不过是等同的位置分配者。中国同志在这方面引用了赫鲁晓夫一个毫不含糊的声明：“我们[苏联和美国]是世界上最强大的国家。如果我们为了和平而联合起来，就不会有战争。如果那时有疯子胆敢发动战争，我们只需要用手指威胁他一下，他就会平静下来。”\footnote{引自《关于战争与和平问题的两条不同路线》，北京，1963年。}

权力本身也变成了一个同质的概念：两个等同的强国，两个超级大国。

联合起来以便每个人都“平静下来”：各就各位。或者反过来（像今天一样）：为了霸权位置而竞争，这种竞争是在对称性中进行的，这种竞争带来了致命、无用、纯粹是位置颠倒的战争：帝国主义间战争，而人民则根据他们的异质力量、根据他们脱离资产阶级平衡计算的力量而起来反抗它，即使这种平衡是恐怖的平衡。

任何被禁锢在不可分割的概念、量化计算和位置分配中的战略，都是一种资产阶级战略，一种帝国主义关于战争与和平的思想。因此，即使被赫鲁晓夫描绘为一种与美国人和平共处和勾结的战略，它也必然反映了准备（帝国主义间）战争的实践。

政治话语中次要过程对主要过程的第一性，因为它让人民扮演了死者的角色，所以它融入了帝国主义屠杀的逻辑。赫鲁晓夫的和平主义学说，就像所有平衡学说一样，是一种侵略性学说：其不可避免的发展遵循着从（表面的）勾结到两个超级大国之间（本质的）竞争的道路。

而这正是中国同志早在1963年就科学地预言的：“[苏共的道路]不是保卫世界和平的道路，而是加剧战争危险并导致战争的道路。”\footnote{《关于战争与和平问题的两条不同路线》，北京，1963年。}（正如我们所看到的，中国共产党第十次代表大会关于苏联和美国之间日益增长的世界大战风险的论点，远非分析的突然改变，而是对1963年预测的、被现实所证实的延续。在十年内，苏联社会帝国主义的内在逻辑——赫鲁晓夫的“和平主义”学说对资产阶级思想隐瞒了它的核心，但对无产阶级辩证法思想揭示了它——展开了其效应。）

相反，无产阶级关于战争与和平的战略，旨在通过发展在主要过程的另一极所累积的所有力量，来中和帝国主义的对称性：这些力量在质上是不同的，即使它们包含并从属于原子武器。

这一战略包括：

\begin{itemize}
    \item 积极支持人民的革命斗争，包括民族解放战争和革命内战（正义战争，主要过程）。然而，在这些战争中，有一个它们的不对称性方面，即人民主要依赖自己的力量（而傀儡和法西斯分子总是主要依赖外国支持）；
    \item 确保核优势，遵循其仅用于防御性用途的质的不同规则，并且从不使用核讹诈，也不指望核武器来实现革命的政治目标。
\end{itemize}

无产阶级战略旨在阻止战争过程的任何致命次要化。

无产阶级思想本质上是主要过程的辩证法。它总是从新事物作为质的肯定而出现的角度来看待事物。矛盾思想必须掌握次要与主要、方面的结构状态与力量的质的动态之间的分裂。它必须同时思考可比较的和不可比较的。它在这里所证明的阶级性质，是它能够在力量关系问题上站队的能力，不是在对称性的空间里，而是在异质性的空间里。

这就是为什么毛主义者会更重视那些从主流思想看来完全不可见或微不足道的工人经验，而不是那些对主流思想来说是决定性的事件。今天在法国，毛主义辩证法家是主要过程的潜水员，他们沉浸在无产阶级的实践深处，在修正主义堆积的次要沉淀物之下。他们所看到和所做的事情，难怪表面上那些无害的对称化者们无法相信自己的眼睛。

真正的辩证问题永远不是摆在第一线的：什么重要的事情正在发生？因为“重要性”这个概念必须根据位置和力量再次分裂。真正的问题永远是：什么新事物正在发生？从这个新事物出发，才需要重新回到重要性上来，就像在人民战争的连续阶段中，一旦红军通过游击战和解放区的缓慢增长，几乎无形地获得了其力量的质的原则，就到了必须在平原上与敌人对抗的时刻，根据敌人自己愚钝的规则，并从量上粉碎他。主要过程的本质规律，不是为了位置的突然互换而进行量的积累，因为那更像是次要过程的对称化道路。

不：主要过程看到未来方面进行其自身的力量积累，这种积累在质上是与它的长期弱点分裂的。在这个积累的终点，最终的问题，那些关于对称性和量的问题，将在它们应有的位置上得到解决：次要的位置。（在革命内战的最后阶段（红军占领中国大城市，越南南方民族解放阵线占领西贡，红色高棉占领金边……），资产阶级思想，作为量化思维的囚徒，对肯定性的差异性视而不见，总是预言着“血腥屠杀”、“绝望的反扑”等等。然而，在所有情况下，这都是最容易、最快速的部分，在这个阶段，傀儡政权与其说是被压垮，不如说是被分解。这是因为质的积累溶解了它们。）

辩证法是关于质的困难和量的便利的理论。位置的互换确实是历史的阶段性标志，但力量的差异化肯定决定了每个阶段的实际内容，即在其本质中，在其分裂中，在其所孕育的下一阶段中，所得到的内容。

\end{document}